% ...........................................................................
% 04c-reencoding.tex -- beats 5-6. The residual re-encoding.
% Legacy anchors sec:methods-vardecomp and sec:methods-holdout are re-\label-ed
% here so v1 cross-references resolve.
% ...........................................................................

\beatsection{R}{5}{The standard target hides the drug}{sec:methods-residual}

\begin{question}[5]
A sentence lists genes in order of expression, so two drugs on the same
cells give almost the same target string and the loss can barely tell them apart.
Can the target be rewritten so that its tokens \emph{are} the drug-specific part?
The instrument counts how many top target positions carry a different gene for
two drugs in one context. Almost none do before the rewrite; a clear majority do
after, at an unchanged retrieval reference.
\end{question}

\noindent
\emph{Encoding} here means how the target string is written, not what the model's
internal states carry (\cref{sec:res-used}). The loss weights every position
equally, so the measurement is a count over the $200$ target positions
(\cref{tab:divergence}).

% ...........................................................................
\begin{table}[htbp]
\begin{threeparttable}
\caption{\textbf{$83\%$ of the standard target's tokens are identical between
any two drugs in the same context.} The residual row is quoted at cell-line
scope; at plate scope it reads $120.3$ and $0.7467$. The middle column is where
the three encodings differ, and it does not on its own identify tokenisation as
the binding constraint.}
\label{tab:divergence}
\small
\begin{tabularx}{\linewidth}{@{}%
  >{\raggedright\arraybackslash}p{0.34\linewidth}
  >{\raggedright\arraybackslash}p{0.28\linewidth}
  >{\raggedright\arraybackslash}X@{}}
\toprule
\thead{Representation} & \thead{Top-$200$ tokens differing between drugs}
  & \thead{Within-well split-half retrieval reference} \\
\midrule
full profile (the standard target) & $34.6$ / $200$ \;($17\%$) & $0.742$ \\
shift (treated $-$ control)        & $111.4$ \;($56\%$)        & $0.742$ \\
\textbf{residual (drug-specific)}  & $118.2$ / $200$ \;($59\%$) & $0.7432$ \\
\bottomrule
\end{tabularx}
\begin{tablenotes}[flushleft]\footnotesize
\item The two columns hold different quantities, a count out of $200$ positions
and a retrieval score, so rows compare within a column and never columns with
each other. The third column pits disjoint cell subsets of one drug, in one cell
line, on one plate, against each other. It is a precision reference, not a
biological replicate ceiling, and it is near-constant across the three rows.
\end{tablenotes}
\end{threeparttable}
\end{table}

\paragraph{The count stops short of the tokenisation claim} The natural
inference, that only a small fraction of the gradient carries drug identity, is
a hypothesis about the optimisation; what is measured is the overlap of target
gene sets, not a token-level loss or a gradient attribution. Gene positions are
not loss-token positions under byte-pair encoding, and the residual target below
alters aggregation, control referencing, generic subtraction, sign coding and
sequence length at once. Separating those five factors needs a matched-budget
factorial this thesis does not run.

\paragraph{The count is not inherited from residual similarity} A descriptive
summary bins $212{,}528$ exhaustive within-group drug pairs on the cosine between
the two residuals. Mean differing top-$200$ genes across the near,
orthogonal and opposite strata run $32.5$, $34.6$ and $35.8$ for the full
profile; $107.4$, $111.3$ and $114.4$ for the shift; $118.4$, $121.0$ and
$118.7$ for the residual, with Spearman coefficients against residual cosine of
$-0.149$, $-0.282$ and $+0.019$. The residual target stays high in every stratum
and lacks the monotone gradient the other two carry, which weakens the reading
that the later swap-distance pattern is inherited trivially from target
divergence without being an independent test of it. The pairs reuse drugs and
groups, so no pair-level inferential claim is made.

% ...........................................................................
\begin{defn}{The drug-specific residual}
For each condition, a (drug, cell line, plate, dose) tuple with at least $40$
treated and $20$ plate-matched control cells,
\begin{equation}
 \resid(d,c) \;=\; \underbrace{\big(\bar{y}_{\textrm{treated}} - \bar{y}_{\textrm{control}}\big)}_{\shift(d,c)}
 \;-\; \underbrace{\frac{1}{|D_c \setminus \{d\}|}\sum_{d' \in D_c,\; d' \neq d} \shift(d',c)}_{\textrm{generic}(c)} .
\end{equation}
The control subtraction removes the cell's own state, leaving the \emph{shift};
the mean over drugs removes the \emph{generic} drug programme every drug
triggers. $D_c$ is the training drugs sharing $c$'s plate and cell line, each
entering once whatever its doses, and the condition's own drug is excluded.
What remains is not what makes this drug different from all drugs. It
is what makes it different \emph{from the other training drugs that happen to
share its plate and cell line}.
\end{defn}

\paragraph{The comparison set is fixed by the plating design} That last sentence
is the price of estimating the generic locally, and mechanism is where it bites.
\Cref{sec:res-scope} finds mechanism-related drugs co-plated $0.404$ of the time
against $0.362$ for a random draw, so for a co-plated mechanism class part of the
class's own shared programme sits inside the generic and is subtracted out of
every member's residual by construction.

That bias was tested. Same-mechanism plate-mates supply a mean $4.4\%$ of a
condition's generic, present for $83\%$ of the $\num{917}$ labelled
mechanism-arm conditions, from Tahoe's \texttt{moa-fine} labels joined to the
plate rosters the channel gate scored. The $\num{155}$ conditions with \emph{no}
same-mechanism plate-mate, whose residuals retain the mechanism programme, give a
gap of \ci{+0.1145}{+0.0276}{+0.2014} against \ci{+0.0737}{+0.0150}{+0.1324} for
the $\num{762}$ that have one, both against the same plate-matched null under the
same two-way clustering. That is the direction the construction predicts, but not
a corrected estimate. The intervals overlap across almost their whole range, the
split is post hoc, and the cleaner stratum rests on $\num{43}$ wells. It supports
a statement about sign, not size.

\begin{scopelimit}[carried into \cref{sec:res-scope} and \cref{sec:limitations}]
Two properties of this frame travel with every number measured in it. First, the
generic is a plate-local mean, so where a mechanism class is co-plated the
mechanism channel is later graded in a frame that has already removed some of
what it is asked to find, which makes an inconclusive verdict on mechanism the
conservative reading and not a generous one. Second, the plate-scoped generic
removes plate structure from the truths and, in expectation, from the
predictions, but that is an expectation and not an identity. An individual
residual may retain some plate structure, and the size of what remains was never
measured. The absolute $\NIR$ values of this frame are read within it and never
set beside the within-plate tables elsewhere in the chapter.
\end{scopelimit}

\paragraph{Which subtraction does the work} Removing the control moves the
between-drug token difference from $34.6$ to $111.4$; removing the generic adds a
further $6.8$ at cell-line scope and $8.9$ at the plate scope this build uses. So
the control subtraction alone buys $92\%$ of the recoverable difference, $90\%$
plate-scoped. The measurement is an overlap of gene sets on one atlas at $K = 200$
($\num{6602}$ conditions, \texttt{target\_divergence.json}), so no margin is
claimed for other atlases or other $K$. It squares with \cref{sec:res-metrics},
where what saturates $\DEdr$ is reversion toward a central point,
\texttt{revert\_center} scoring $0.9999$ without fitting anything.

The generic is not dispensable. It carries $39\%$ of the shift's squared norm
(\cref{sec:res-transfer}), and subtracting it isolates the drug-specific
remainder once the baseline is gone. The inclusion thresholds sit where the
identifiability curve of \cref{sec:res-blind} flattens, the within-well
split-half precision reference being $0.61$ at $10$ cells, $0.76$ at $40$ and
$0.85$ at $120$.

% ...........................................................................
\paragraph{A definition that could have been circular, and what breaks it} Four
properties of how the generic is estimated decide whether the remainder can be
trusted. First, the shape of the problem. A condition enters training only if its
residual reproduces across a half-split; a residual exists only once the generic
has been subtracted; the generic is fitted on the training set; and the training
set is what the reliability line selects. That is a loop
(\cref{fig:circularity}, left), and an earlier build closed it, with generic and
filter estimated over every condition and the holdout assigned afterwards, so
each held-out target was defined partly by other held-out conditions.

The loop is cut by ordering. Inventory and split read metadata only, before any
pseudobulk exists; fit hands the generic the training keys explicitly; transform
projects held-out conditions through the fitted generic, which they never enter
(\cref{fig:circularity}, right). Because the split is made on metadata alone,
nothing downstream of it can have defined it. A poison test makes the ordering
checkable in both directions.\gloss{poison test}{Corrupt the held-out side; the
training targets must not move. Corrupt a training condition; they must.} The
other three properties fix choices the staging leaves open.

% ===========================================================================
% circularity.tex -- the residual definition's loop, and the ordering that cuts
% it. \input by Sections/04c-reencoding.tex at sec:methods-residual. Compiles
% against thesis_v2/preamble.tex and nothing else: no external data, no
% \includegraphics, no package this document does not already load
% (tikz + arrows.meta, both in preamble.tex section 9).
%
% THIS IS A PROVENANCE DIAGRAM, NOT A RESULT. Deliberately absent: every NIR
% value, every arm, every split-level number. Two counts appear, once, and both
% are already in the prose of this same section. Sibling circularity.json lists
% every number drawn.
%
% THE LEFT PANEL CARRIES NO BUILD LABEL, DELIBERATELY. It drew "THE EARLIER
% BUILD -- NOT THE SHIPPED ONE" until this revision, which spent half a
% full-width figure on a pipeline the thesis did not run. The historical fact
% -- that an earlier build did close this loop -- is a disclosure, and a
% disclosure belongs in prose where it can be qualified: it is stated at
% Sections/04c-reencoding.tex:148-151 ("That is a loop [...], and an earlier
% build closed it, with generic and filter estimated over every condition and
% the holdout assigned afterwards"). That sentence stands on its own and was
% NOT deleted when the label came off the panel. What the left panel draws now
% is the structural hazard itself, in the present tense and attached to no
% build: the cycle the four definitions form, and the edge that fitting before
% splitting adds to it. The two panels therefore read hazard and remedy, which
% is the only reason the staged ordering is load-bearing rather than fastidious.
%
%   2,523 of 4,999 training conditions retained (50.5%)   04c-reencoding.tex:211-212
%   the resolved reliability line, +0.1086                 04c-reencoding.tex:211
%
% Cross-checked, all four agreeing, before drawing:
%   Sections/03-Apparatus.tex:206     "2,523 of 4,999 train (50.5\%)"
%   Sections/04c-reencoding.tex:212   "$2{,}523$ of $\num{4999}$ ... ($50.5\%$)"
%   Sections/04d-baselines.tex:337    "only $\num{2523}$ are training"
%   Sections/05-Limitations.tex:33    "$\num{2523}$ of $\num{4999}$ ... $\cos > 0.109$"
% +0.109 is the measured null's 95th percentile and +0.1086 is what it resolves
% to; the prose prints both, and the node carries the resolved value because
% that is the one the builder applies.
%
% LAYOUT NOTES, all forced by a build:
%   * fullwidth = 174mm = 493.23pt (textwidth 142mm + marginparsep 4mm +
%     marginparwidth 28mm, preamble.tex:88-90 and :525). The bounding box is
%     PINNED to 17.35cm = 491.81pt, 1.4pt inside the measure. Unpinned, the
%     picture measured 3.10pt OVER and the figure body reported an overfull
%     hbox: nodes with anchor=north west hang their 3.33pt inner sep left of
%     the anchor at x=0.20, and the strip rules end at x=17.35. Pinning clips
%     no ink -- no drawn object lies outside [0, 17.35] x [-0.95, 7.95].
%   * NO node uses text width. Every line break is written by hand. The first
%     build did use text width and TeX hyphenated inside the boxes
%     ("the relia-/bility line"), grew four of them past their minimum heights
%     and overlapped the rows below. Hand-broken lines make every box width and
%     height deterministic, which is what lets the coordinates below be trusted.
%   * the four objects are ONE node set drawn twice, same style, same wording,
%     same 2.80cm minimum width, so the right panel reads as a re-laying of the
%     left and not as a second diagram. Only the arrangement changes.
%   * the fatal edge is drawn and STRUCK, not deleted. A deleted edge leaves
%     nothing for the caption to point at, and the passage is about which single
%     edge the ordering removes.
%   * the two streams leaving SPLIT must not cross. They cannot both turn upward
%     in the order they leave: the stream that leaves HIGHER (held-out, 2.2mm
%     above the east anchor) turns up FIRST (x=10.30) and the stream that leaves
%     LOWER (training keys, at the east anchor) turns up LATER (x=10.85). Any
%     other pairing crosses, and a crossing here reads as contact between them.
%   * both streams use -| rather than an explicit corner coordinate, so the
%     corner inherits the node anchor's own y. Written as a literal corner, the
%     first segment came out as a visible diagonal whenever a node's height
%     differed by a hair from the value assumed at design time.
%   * the held-out stream is routed over the top of the fit column, and its one
%     arrowhead points right, into TRANSFORM. No arrowhead anywhere in the right
%     panel points back left or up into the fit; that absence is the claim.
%   * NO edge label is masked onto its arrow. The first build masked all of
%     them and the labels ate the edges: the C->D gap in the left panel is
%     1.20cm and "subtracted to give" is 1.15cm wide, so the bottom edge of the
%     cycle disappeared and the loop stopped reading as a loop. Every label now
%     sits beside its edge. The tightest is "training keys", hung under the
%     corner at x = 10.43 to 11.27, between the held-out riser at 10.30 and the
%     fit column's left edge at 11.35.
% ===========================================================================
\begin{figure}[htbp]
\begin{fullwidth}
\centering
\begin{tikzpicture}[x=1cm, y=1cm]
  \useasboundingbox (0,-0.95) rectangle (17.35,7.95);
  \tikzset{
    % the shared node set. Identical style in both panels.
    obj/.style={draw=rulegrey, line width=0.4pt, fill=white, inner sep=3pt,
                align=center, font=\scriptsize, text=ink,
                minimum width=2.80cm, minimum height=0.68cm},
    edg/.style={-{Stealth[length=3.6pt,width=2.8pt]}, draw=ink,
                line width=0.5pt},
    lab/.style={font=\scriptsize, text=quietink, fill=white, inner sep=1pt,
                align=center},
    side/.style={font=\scriptsize, text=quietink, align=flush right,
                 anchor=east, text width=3.20cm},
    quiet/.style={font=\scriptsize\itshape, text=quietink, align=center},
    tag/.style={font=\sffamily\scriptsize\bfseries, text=accent},
    sub/.style={font=\footnotesize\itshape, text=quietink}}
  \def\T#1{{\sffamily\scriptsize\bfseries\color{accent}#1}\\[1.0pt]}
  \def\Q#1{\\[1.0pt]{\color{quietink}#1}}

  % ======================= PANEL HEADS =====================================
  % Hazard and remedy, one head each, same shape: a tag naming the structure
  % and a sub saying what happens to the four objects. Neither names a build.
  \node[tag, anchor=north west] at (0.20,7.90) {THE CIRCULARITY};
  \node[sub, anchor=north west] at (0.20,7.55)
    {the four objects define one another};
  \node[tag, anchor=north west] at (7.50,7.90) {THE ORDERING THAT CUTS IT};
  \node[sub, anchor=north west] at (7.50,7.55)
    {the same four objects, staged by execution order};

  \draw[rulegrey, line width=0.4pt] (7.20,1.15) -- (7.20,6.60);

  % ======================= LEFT: THE CLOSED LOOP ===========================
  \node[obj, anchor=north] (A) at (1.60,6.10) {the reliability line};
  \node[obj, anchor=north] (B) at (5.60,6.10) {the training set};
  \node[obj, anchor=north] (C) at (5.60,4.45) {the generic};
  \node[obj, anchor=north] (D) at (1.60,4.45) {the residual};

  % NONE of the four loop labels is masked onto its arrow. Masked, they ate the
  % edges: the C->D gap is 1.20cm and "subtracted to give" is 1.15cm wide, so
  % the bottom edge of the cycle vanished entirely and the loop stopped reading
  % as a loop. Each label now sits beside its own edge -- the two horizontals
  % labelled outside the cycle, the two verticals inside it -- and all four
  % edges draw unbroken.
  \draw[edg] (A.east) -- (B.west);
  \draw[edg] (B.south) -- (C.north);
  \draw[edg] (C.west) -- (D.east);
  \draw[edg] (D.north) -- (A.south);
  \node[lab, anchor=south] at (3.60,5.86) {selects};
  \node[lab, anchor=east]  at (5.45,4.94) {fits};
  % clears the box row's bottom edge at y=3.77: at y=3.97 its first line sat
  % level with the two boxes' lower corners, 0.02cm from each, and read as a
  % collision even though it was not one.
  \node[lab, anchor=north] at (3.60,3.88) {subtracted\\to give};
  \node[lab, anchor=west]  at (1.75,4.94) {must reproduce\\across a half-split};

  % --- the fatal edge: drawn once, struck once, and labelled --------------
  % Present tense, and attached to an ordering rather than to a build: this is
  % what fitting before splitting does to the cycle, not what one run once did.
  % The quiet italic under H answers "why would they enter the fit?" inside the
  % panel, and mirrors the right panel's one quiet italic under SPLIT, so each
  % half carries exactly one such annotation.
  \node[obj, anchor=north] (H) at (1.60,2.60) {held-out conditions};
  \draw[edg] (H.east) -| (C.south);
  \node[lab, anchor=east] at (5.45,2.95) {enter the fit};
  % Its y is the strike caption's y, not a value of its own: both are two-line
  % blocks of the same size, so sharing anchor y = 1.74 gives them a common top
  % AND a common bottom (0.944) and the bottom-left reads as one row rather than
  % two blocks at two heights. Widest line measures 54.42pt = 1.91cm, so with
  % inner sep the node spans x 0.53-2.67, clear of the strike caption's left
  % edge at x = 3.10 by 0.43cm. Top sits 0.18cm under H, matching the 0.20cm
  % the right panel leaves under SPLIT.
  \node[quiet, anchor=north] at (1.60,1.74)
    {when the fit\\precedes the split};
  % the strike. One stroke. Nothing else in this figure is accent-coloured
  % except the structural tags.
  \draw[accent, line width=1.1pt] (4.25,1.94) -- (4.75,2.58);
  \node[font=\sffamily\scriptsize, text=accent, anchor=north west,
        align=left] at (3.10,1.74)
    {struck: the one edge\\the ordering removes};

  % ======================= RIGHT: THE STAGED PIPELINE ======================
  \fill[panelbg] (7.45,3.55) rectangle (10.15,6.30);
  \node[obj, anchor=north, minimum width=2.55cm] (INV) at (8.80,6.10)
    {\T{INVENTORY}wells, drugs, doses};
  \node[obj, anchor=north, minimum width=2.55cm] (SPL) at (8.80,5.10)
    {\T{SPLIT}train \textbar{} held out};
  \draw[edg] (INV.south) -- (SPL.north);
  \node[quiet, anchor=north] at (8.80,4.22)
    {metadata only:\\no expression data};

  % FIT, and the training side that descends from it
  \node[obj, anchor=north] (G) at (12.75,6.10) {\T{FIT}the generic};
  \node[obj, anchor=north] (R) at (12.75,4.60) {the residual};
  \node[obj, anchor=north] (L) at (12.75,3.30)
    {the reliability line\Q{$+0.1086$}};
  \node[obj, anchor=north] (S) at (12.75,1.90) {the training set};
  \draw[edg] (G.south) -- (R.north);
  \draw[edg] (R.south) -- (L.north);
  \draw[edg] (L.south) -- (S.north);

  \node[side] at (11.25,2.80) {$2{,}523$ of $4{,}999$ training conditions
    retained ($50.5\%$)};
  \node[side] at (11.25,1.56) {the residual targets the model is trained on};

  % training keys, handed to the fit explicitly
  % -| , not an explicit corner coordinate: the corner must inherit SPL.east's
  % own y, or the first segment comes out as a visible diagonal when the node's
  % height differs by a hair from the value assumed here.
  % This is the one edge label NOT masked onto its arrow. The riser is 1.0cm
  % long and a masked label ate 0.6cm of it, leaving the arrow looking like two
  % disconnected stubs. It hangs under the corner instead, in the only clear
  % space in the panel: x 10.43-11.27, below the horizontal and above the
  % retention annotation, which tops out at y = 3.25.
  \draw[edg] (SPL.east) -| (10.85,5.675) -- (G.west);
  \node[font=\scriptsize, text=quietink, align=center, anchor=north]
    at (10.85,4.50) {training\\keys};

  % TRANSFORM, fed by the fitted generic and by a stream that bypasses the fit
  \node[obj, anchor=north] (TR) at (15.95,6.10)
    {\T{TRANSFORM}held-out conditions\\projected through\\the fitted generic};
  \draw[edg] (G.east) -- (G.east -| TR.west);
  \draw[edg] ([yshift=2.2mm]SPL.east) -| (10.30,6.45) -- (15.95,6.45)
    -- (TR.north);
  \node[font=\scriptsize, text=quietink, anchor=south] at (13.10,6.55)
    {held-out keys --- they never enter the fit};

  \node[obj, anchor=north] (HT) at (15.95,4.30)
    {held-out residual\\targets};
  \draw[edg] (TR.south) -- (HT.north);
  \node[quiet, anchor=north] at (15.95,3.25)
    {deliberately unfiltered};
  \node[quiet, anchor=north] at (15.95,2.80)
    {nothing here\\re-enters the fit};

  % ======================= THE POISON TEST STRIP ===========================
  \draw[rulegrey, line width=0.4pt] (0.20,0.85) -- (17.35,0.85);
  \node[tag, anchor=north west, align=left] at (0.20,0.46)
    {THE POISON\\TEST};
  \node[font=\sffamily\scriptsize\bfseries, text=ink, anchor=north west]
    at (3.30,0.55) {CORRUPTION APPLIED};
  \node[font=\sffamily\scriptsize\bfseries, text=ink, anchor=north west]
    at (9.60,0.55) {REQUIRED OUTCOME};
  \draw[rulegrey, line width=0.4pt] (3.30,0.20) -- (17.35,0.20);
  \node[font=\sffamily\footnotesize, text=ink, anchor=north west]
    at (3.30,0.08) {corrupt every held-out expression vector};
  \node[font=\sffamily\footnotesize, text=ink, anchor=north west]
    at (9.60,0.08) {the training targets are byte-identical};
  \node[font=\sffamily\footnotesize, text=ink, anchor=north west]
    at (3.30,-0.40) {corrupt one training condition};
  \node[font=\sffamily\footnotesize, text=ink, anchor=north west]
    at (9.60,-0.40) {the training targets must move};
  \draw[rulegrey, line width=0.4pt] (3.30,-0.85) -- (17.35,-0.85);
\end{tikzpicture}
\caption[The residual's circular definition, cut by
ordering]{\textbf{The definition of the residual could have closed on itself,
and what stops it is execution order, not an argument.} \emph{Left, the
circularity.} A condition enters training only if
its residual reproduces across a half-split; the residual exists only once the
generic is subtracted; the generic is fitted on the training set; and the
training set is what the reliability line selects. Estimate the generic and the
filter over every condition and assign the holdout afterwards, and held-out
conditions enter the fit --- the struck edge --- so each held-out
target is defined partly by other held-out conditions. \emph{Right, the
ordering that cuts it,} the build this thesis reports. The same four objects,
staged. Inventory and split read
metadata only, before any pseudobulk exists, so nothing downstream of the split
can have defined it; the fit is handed the training keys explicitly; and
transform projects held-out conditions through the fitted generic in one
direction, with no path back into it. The generic is the mean over training
drugs other than this one, within plate. The two counts are the
\emph{training-set} retention at the resolved reliability line $+0.1086$, the
$95$th percentile of a null pairing half $A$ of one condition with half $B$ of
another; the held-out sets are deliberately unfiltered, since dropping held-out
conditions for failing their own reproducibility bar would select the
evaluation set on its own outcome. \emph{Beneath,} the two assertions that make
the ordering checkable in both directions rather than merely asserted. This is
a provenance diagram: no $\NIR$ value, arm or split-level result appears in
it.}
\label{fig:circularity}
\end{fullwidth}
\end{figure}


\paragraph{Second, the scope is the plate} The original argument for cell-line
scope was reproducibility. $61.7\%$ of conditions reproduce across a half-split
at cell-line scope against $19.3\%$ at plate scope (cosine threshold $0.2$,
$n = 6602$), on the reasoning that a plate holds too few drugs to estimate a
mean.\corrmark{the retention argument for cell-line scope is
retracted here; the scope decision now rests on the two diagnostics below} Both
figures were computed with a shared control pseudobulk subtracted from both
halves, which makes the halves correlate through their common control error, and
the inflation favours cell-line scope precisely because the cell-line generic is
itself less noisy. That alone retracts the retention argument. A disjoint-control
recomputation exists but was never written to an artifact, and the two surviving
records disagree on the plate figure, so no corrected percentages are quoted for
this $n = 6602$ comparison. Both agree that plate scope retains at least as many
conditions, so retention is no reason to prefer cell-line scope.

Two diagnostics settle the scope in plate's favour, both reported in
\cref{sec:res-transfer}. At cell-line scope the dose-transfer ordering inverts,
and the shared-minus-split control bias reads $+0.242$ against $-0.000$ at plate
scope. A plate group with fewer than three other training drugs yields no
generic, and the condition is dropped and counted.

\paragraph{Third, the mean is over drugs and leaves the condition's own drug out}
Two things force that. A drug measured at five doses would otherwise carry five
times the weight of a single-dose drug; and once the generic is train-only, a
training condition's generic would contain its own drug while a held-out
condition's would not, so the two splits would be scored against differently
defined targets. The generic is therefore the mean over training drugs other
than this one, within scope.

\paragraph{Fourth, the reliability line is resolved from the data} The threshold
was the largest single determinant of the training set, and it had been
inherited. The inherited $0.2$ came from an era when reliability was measured
with shared controls, a different quantity on a different scale, under which it
retained the $61.7\%$ above against the $19.0\%$ ($\num{948}$ of the $\num{4999}$
training candidates) it retains once each half is measured against its own
control half. The two are not a matched pair, differing in population and scope
as well as control handling, but between them a moderate bar had silently become
a stringent one.

\begin{defn}{The reliability line}
A condition is kept for training only if its residual reproduces across a
half-split, $\cos(\resid_A, \resid_B)$ above a threshold, with each half
measured against its own control half so that the two share nothing. No new
threshold is picked by hand. The builder reports the whole retention curve
against a null that pairs half $A$ of one condition with half $B$ of a
\emph{different} condition, same encoding, same dimensionality, same noise, no
shared biology, and sets the line at that null's $95$th percentile. A condition
is kept when its two halves agree more than two unrelated conditions' halves do.
At the measured null ($95$th percentile $+0.109$, resolving to $+0.1086$) this
retains $2{,}523$ of $\num{4999}$ training conditions ($50.5\%$).
\end{defn}

The inherited $0.2$ survives as an argparse default in three downstream configs,
\texttt{truth\_repro\_thr} in \texttt{re\_v3.json} and \texttt{repro\_thr} in
\texttt{vardecomp\_matched.json} and \texttt{channel\_gate\_v3.json}. It did not
act as the training filter. Of the $\num{300}$ scored training conditions,
$\num{181}$ sit below $0.2$, the minimum is $+0.109$, and none falls at or below
the resolved $+0.1086$. The held-out sets are deliberately unfiltered, since
dropping held-out conditions for failing their own reproducibility bar would
select the evaluation set on its own outcome; noise in retained training
conditions costs learning efficiency, not validity.

The half-split agreement is a within-well sampling-precision criterion. It is not
biological reproducibility, since $\num{286}$ of $\num{287}$ distinct (drug,
dose) treatments in this cache occur in exactly one well. The retained fraction
says how much drug-specific signal survives pseudobulk noise at this cell budget,
not how much biology is real.

\paragraph{How the residual becomes a string, and how it is scored}
\Cref{fig:frame} draws the path, from the two subtractions to the vector that is
scored. The target is \texttt{<up genes> \DOWN{} <down genes> \ENDCELL}, $100$
genes per block, ordered by signed residual magnitude; sign is biology, so it
gets a symbol.
Predicted and true residuals are both mapped to a signed rank vector, top-$100$
up positive, top-$100$ down negative, weight $1/\log_2(\textrm{rank}+1)$;
similarity is cosine, for the reason given in \cref{sec:res-metrics}, where the
residual row of \cref{tab:sigma} declares it as this frame's $\sigma$.
A generation supplies its ranks by token order within each
block, and a token outside the panel or already placed is skipped without
consuming a rank, so a block fills to at most $100$ genes. The sampled
generations of one condition are averaged as vectors and the mean is scored
once, not scored one by one and averaged.
The score is $\NIR$ over other drugs in the same cell line. That set
spans plates, in apparent violation of the same-plate rule of
\cref{sec:res-blind}, and stands only because the plate-scoped generic has
already subtracted each plate's mean shift, in the expectation sense fenced
off above.

% ===========================================================================
% fig-frame.tex -- v2. The residual frame, the encoding fork, and the SCORED
% object. \input by Sections/04c-reencoding.tex. Compiles against
% thesis_v2/preamble.tex and nothing else; no external data, no
% \includegraphics. FULLWIDTH, and MEASURED rather than assumed: the
% tikzpicture's own bounding box is 491.31pt wide against a fullwidth measure
% of 495.08pt (\textwidth 404.03pt + \marginparwidth + \marginparsep). The
% measurement is made by the picture itself, guarded by \fframedebug;
% figs/_harness_frame.tex switches it on, and it is inert in main.tex.
%
% RULE INHERITED FROM preamble.tex section 8: a page carrying a fullwidth float
% carries NO \marginpar. Every note in this figure is therefore drawn INSIDE
% the tikzpicture. Nothing here calls \gloss, \seealso or any margin wrapper.
%
% EVERY NUMBER IN THIS FIGURE, WITH ITS SOURCE. Nothing is invented.
%   34.6 / 200  (17%)  full profile, cell-line scope
%   118.2 / 200 (59%)  residual,     cell-line scope
%   120.3              residual,     plate scope            -- caption only
%   0.742 -> 0.743     within-well split-half retrieval reference, full ->
%                      residual (stored 0.741986901358855 and 0.7432424792)
%     RESULTS_cluster/target_divergence.json, scopes.cellline.{full,residual}
%     .divergence_mean and .ceiling_retrieval, and scopes.plate.residual;
%     and tab:divergence, Sections/04c-reencoding.tex lines 42-44.
%   100 genes per block, [DOWN], [END_CELL], weight 1/log2(rank+1), cosine,
%   NIR over other drugs in the same cell line
%     Sections/04c-reencoding.tex, "How the residual becomes a string, and how
%     it is scored"; and endcell/analysis/residual_eval.py,
%     signed_rank_from_sentence / signed_rank_from_vector (k = 100,
%     v[gi] = sgn / np.log2(r + 1.0)).
%   946 panel genes -- Sections/04c-reencoding.tex, "a 946-dimensional profile".
%   the stem weights: 1/log2(r+1) for r = 1..8 =
%     1.0000, 0.6309, 0.5000, 0.4307, 0.3869, 0.3562, 0.3333, 0.3155.
%     Every stem in stage 3 is one of these eight numbers; the +1 / -1 ticks
%     are there so a reader can check them.
%   "fewer than three other training drugs yields no generic and the condition
%     is dropped" -- Sections/04c-reencoding.tex, "Second, the scope is the
%     plate".
%   [DOWN] registered atomic / unregistered splits into sub-words --
%     Sections/03-Apparatus.tex line 283, and 00-HowToRead.tex line 209.
%
% WHAT IS ILLUSTRATIVE, AND SAYS SO IN THE CAPTION:
%   * the six gene profiles in stage 1. No axis is drawn; they carry no
%     measurement. Their ARITHMETIC is exact bar by bar, which is the only
%     thing a reader may take from them:
%       treated  {17, 9,14,12, 8, 6,11, 5}
%       control  {10,10, 8, 8, 7, 7,10, 6}
%       shift    { 7,-1, 6, 4, 1,-1, 1,-1}  = treated - control
%       generic  { 3, 2, 3, 2,-2, 2,-2, 2}
%       residual { 4,-3, 3, 2, 3,-3, 3,-3}  = shift - generic
%     Chosen so the residual carries the MAJORITY of the shift's squared norm
%     (74 of 106 here), because fig:residual-energy measures 62.1%
%     drug-specific and a figure drawing the residual as a sliver would
%     contradict it. Both figures are UNCENTRED sums of squares, each a
%     squared L2 norm and NOT a variance: the shift {7,-1,6,4,1,-1,1,-1} has
%     mean 2, not 0, and 49+1+36+16+1+1+1+1 = 106; the residual gives 74 the
%     same way. The label fig:residual-energy and the file
%     figs/04c-residual-energy.tex keep the older spelling because they are a
%     cross-reference target and an \input target.
%     v1 drew {1,0,0,1,0,1,0,0}: five of its eight bars had height zero and
%     rendered as nothing at all, so the one object the thesis is about read
%     as "missing" or "to be determined". FIXED HERE, and drawn SIGNED,
%     because the sign is what splits the two blocks in stage 2.
%   * which gene sits at which coordinate in stage 3, and which two the
%     prediction puts on the wrong side of the boundary. The STEM HEIGHTS are
%     not illustrative -- they are the exact weights.
%   * no cosine VALUE is printed anywhere. The two stem vectors are drawn, not
%     measured, so a number there would be an invented result.
%
% LAYOUT NOTES, all forced by a build:
%   * three stages, two dividers, one header rule, one brace. The brace spans
%     all THREE because the retrieval reference is an end-of-pipeline quantity:
%     it is a retrieval score computed in the scored space of stage 3, so
%     stopping it at stage 2 would misattribute it.
%   * stage 2's target string is ONE line. Wrapped over two it read as two
%     separate sequences and lost the up -> [DOWN] -> down ordering, which is
%     the only reason [DOWN] is drawn at all. Every chip is a NAMED node
%     chained off its predecessor's east anchor, so the two braces attach to
%     real node boundaries rather than to hand-computed x's. That was not
%     fastidiousness: the hand estimate of the string's width was 6mm short of
%     the truth, and on the first build [END_CELL] crossed the divider into
%     stage 3.
%   * type sizes follow figs/circularity.tex: \footnotesize (10pt, the caption
%     size) for the sentences a reader must read, \scriptsize (8pt) for symbol
%     labels, chips, brace labels and subordinate notes. Nothing is \tiny.
%   * stage 1 labels are centred on profiles 1.85cm apart. At 1.45cm
%     "generic(c)" and "residual r(d,c)" touch (v1's note, still true). They
%     are \scriptsize and not \footnotesize for the same reason: at 10pt they
%     collide even at 1.85cm, which a build showed.
% ===========================================================================
\begin{figure}[htbp]
\begin{fullwidth}
\centering
\begin{tikzpicture}[
  x=1cm, y=1cm, font=\small,
  bar/.style={fill=ink!72, draw=none},
  faintbar/.style={fill=ink!30, draw=none},
  base/.style={rulegrey, line width=0.25pt},
  gene/.style={draw=rulegrey, rounded corners=1pt, inner sep=1.6pt,
               font=\ttfamily\scriptsize, text=ink, anchor=west},
  shared/.style={gene, fill=rulegrey!12, text=quietink, draw=rulegrey!55},
  differs/.style={gene, fill=rulegrey!32, draw=ink!75, thick},
  atomic/.style={rounded corners=1pt, inner sep=1.6pt, fill=ink, draw=ink,
                 font=\ttfamily\scriptsize, text=white, anchor=west},
  dots/.style={anchor=west, font=\scriptsize, text=quietink},
  lbl/.style={font=\scriptsize, text=quietink},
  strong/.style={font=\scriptsize, text=ink},
  said/.style={anchor=west, font=\footnotesize, text=ink},
  op/.style={font=\normalsize, text=quietink},
  stagehead/.style={anchor=west, font=\sffamily\footnotesize\bfseries,
                    text=accent},
  note/.style={anchor=north west, font=\scriptsize, text=quietink,
               align=flush left, inner sep=0pt},
]

% ---- one bar profile. #1 x, #2 baseline y, #3 signed heights, #4 style.
% Heights are in "units"; one unit is 0.052cm. Bars are 0.125cm wide on a
% 0.160cm pitch, so eight span 1.245cm. The baseline rule is drawn under EVERY
% profile, because three of the six are signed and a bar below the line has to
% read as negative rather than as a misprint.
\newcommand{\fframeprof}[4]{%
  \foreach \hh [count=\ii from 0] in {#3} {%
    \pgfmathsetmacro{\bxx}{#1 + \ii*0.160}%
    \pgfmathsetmacro{\byy}{\hh*0.052}%
    \fill[#4] (\bxx,#2) rectangle ++(0.125,\byy);%
  }%
  \draw[base] (#1,#2) -- ++(1.245,0);%
}

% ---- one signed-rank stem plot. #1 zero-line y, #2 slot/weight list.
% Slots sit on a 0.255cm pitch from x = 13.16; a weight of 1 is 0.34cm.
\newcommand{\fframestems}[2]{%
  \draw[base] (13.10,#1) -- (17.02,#1);%
  \foreach \ss/\ww in {#2} {%
    \pgfmathsetmacro{\sxx}{13.16 + (\ss-1)*0.255}%
    \pgfmathsetmacro{\syy}{#1 + \ww*0.34}%
    \draw[ink!72, line width=0.55pt] (\sxx,#1) -- (\sxx,\syy);%
    \fill[ink!72] (\sxx,\syy) circle (0.5pt);%
  }%
  \draw[rulegrey, line width=0.3pt] (12.96,{#1-0.34}) -- (12.96,{#1+0.34});%
  \draw[rulegrey, line width=0.3pt] (12.92,{#1+0.34}) -- (13.00,{#1+0.34});%
  \draw[rulegrey, line width=0.3pt] (12.92,{#1-0.34}) -- (13.00,{#1-0.34});%
  \node[anchor=east, font=\scriptsize, text=quietink]
       at (12.88,{#1+0.34}) {$+1$};%
  \node[anchor=east, font=\scriptsize, text=quietink]
       at (12.88,{#1-0.34}) {$-1$};%
}

% ================================================== the header band
\node[stagehead] at (0.05,3.80) {1\,\textperiodcentered\,the residual is built};
\node[stagehead] at (5.70,3.80) {2\,\textperiodcentered\,it becomes a string};
\node[stagehead] at (12.60,3.80) {3\,\textperiodcentered\,it becomes a vector};
\draw[rulegrey, line width=0.4pt] (0.05,3.62) -- (17.30,3.62);
\draw[rulegrey!70, line width=0.3pt] (5.50,-1.66) -- (5.50,3.54);
\draw[rulegrey!70, line width=0.3pt] (12.40,-1.66) -- (12.40,3.54);

% ================================================== STAGE 1: two subtractions
\node[lbl]    at (0.67,3.36) {treated};
\node[lbl]    at (2.52,3.36) {control};
\node[strong] at (4.37,3.36) {shift $\shift(d,c)$};
\fframeprof{0.05}{2.30}{17,9,14,12,8,6,11,5}{bar}
\node[op] at (1.53,2.58) {$-$};
\fframeprof{1.90}{2.30}{10,10,8,8,7,7,10,6}{bar}
\node[op] at (3.38,2.58) {$=$};
\fframeprof{3.75}{2.30}{7,-1,6,4,1,-1,1,-1}{bar}

\node[strong] at (0.67,1.90) {shift};
\node[lbl]    at (2.52,1.90) {generic$(c)$};
\node[strong] at (4.37,1.90) {residual $\resid(d,c)$};
\fframeprof{0.05}{1.30}{7,-1,6,4,1,-1,1,-1}{bar}
\node[op] at (1.53,1.42) {$-$};
\fframeprof{1.90}{1.30}{3,2,3,2,-2,2,-2,2}{faintbar}
\node[op] at (3.38,1.42) {$=$};
\fframeprof{3.75}{1.30}{4,-3,3,2,3,-3,3,-3}{bar}

\node[note, text width=5.30cm] (nG) at (0.05,0.88)
  {generic$(c)$ is the mean over the \emph{other} training drugs in the same
   plate group; a group with fewer than three of them yields no generic and the
   condition is dropped. The residual is drawn signed because its sign is what
   splits the two blocks in stage~2.};

% ================================================== STAGE 2: the encoding fork
\node[said] at (5.62,3.30) {\textbf{(a)} genes ranked by \emph{expression}};
\node[shared]  (a1) at (5.62,2.88) {ACTB};
\node[shared]  (a2) at ([xshift=3pt]a1.east) {GAPDH};
\node[shared]  (a3) at ([xshift=3pt]a2.east) {TPT1};
\node[differs] (a4) at ([xshift=3pt]a3.east) {CDKN1A};
\node[shared]  (a5) at ([xshift=3pt]a4.east) {EEF1A1};
\node[dots]    (a6) at ([xshift=4pt]a5.east) {$\cdots$};
\node[shared]  (b1) at (5.62,2.48) {ACTB};
\node[shared]  (b2) at ([xshift=3pt]b1.east) {GAPDH};
\node[shared]  (b3) at ([xshift=3pt]b2.east) {TPT1};
\node[differs] (b4) at ([xshift=3pt]b3.east) {HMOX1};
\node[shared]  (b5) at ([xshift=3pt]b4.east) {EEF1A1};
\node[dots]    (b6) at ([xshift=4pt]b5.east) {$\cdots$};
\node[said] at (5.62,2.06)
  {only \textbf{34.6 / 200} tokens differ \;(\textbf{17\%})};

\node[said] at (5.62,1.56)
  {\textbf{(b)} genes ranked by \emph{signed residual}};
\node[differs] (c1) at (5.62,1.14) {CDKN1A};
\node[differs] (c2) at ([xshift=3pt]c1.east) {MDM2};
\node[dots]    (c3) at ([xshift=3pt]c2.east) {$\cdots$};
\node[atomic]  (c4) at ([xshift=4pt]c3.east) {[DOWN]};
\node[differs] (c5) at ([xshift=4pt]c4.east) {HMOX1};
\node[dots]    (c6) at ([xshift=3pt]c5.east) {$\cdots$};
\node[atomic]  (c7) at ([xshift=4pt]c6.east) {[END\_CELL]};
\draw[decorate, decoration={brace, mirror, amplitude=3pt}, rulegrey]
  (c1.south west) ++(0,-0.06) -- ($(c3.south east)+(0,-0.06)$);
\node[lbl, anchor=north west] at ($(c1.south west)+(0,-0.16)$)
  {$100$ genes, $\resid > 0$};
\draw[decorate, decoration={brace, mirror, amplitude=3pt}, rulegrey]
  (c5.south west) ++(0,-0.06) -- ($(c6.south east)+(0,-0.06)$);
\node[lbl, anchor=north west] at ($(c5.south west)+(0,-0.16)$)
  {$100$ genes, $\resid < 0$};
\node[said] at (5.62,0.20)
  {\textbf{118.2 / 200} tokens differ \;(\textbf{59\%})};

\node[atomic] (dn) at (5.62,-0.28) {[DOWN]};
\node[note, text width=5.35cm] (nD) at ($(dn.east)+(0.14,0.11)$)
  {is \emph{one} registered atomic token. Unregistered it splits into
   sub-words, and the up\slash down boundary stops being a clean symbol.};

% ================================================== STAGE 3: the scored vector
\node[note, text width=4.70cm] (nW) at (12.60,3.44)
  {signed rank weights $w_r = 1/\log_2(r{+}1)$};

% the two coordinates on which the prediction and the truth carry OPPOSITE
% SIGNS -- the boundary error the [DOWN] token exists to prevent. Drawn FIRST
% so the stems sit on top of it.
\fill[rulegrey!18] (14.82,0.70) rectangle (15.33,2.42);
% The band's label sits in the GAP BETWEEN the two plots, running right from the
% band's right edge. Above the plots it ran into the second line of the weights
% note, which is math and taller than a text line; centred on the band it set
% flush against "true v" and the two read as one phrase. In the gap the nearest
% ink is the predicted plot's slot-10 stem, which stops at y = 1.72.
% anchor=EAST at 17.30, not west at 15.42: anchored west, this one 2.07cm label
% was the widest thing in the picture and pushed the bounding box out to
% 497.60pt, 1.35pt past the fullwidth measure -- one Overfull \hbox, which
% tools/build.sh caps at zero.
\node[lbl, anchor=east] at (17.30,1.46) {opposite signs};

\node[strong, anchor=west] at (13.16,2.50) {predicted $\hat{v}$};
\fframestems{2.06}{1/0.6309, 2/1.0000, 3/0.5000, 4/0.3869, 5/0.4307,
                   6/0.3562, 7/0.3155, 8/-0.6309, 9/0.3333, 10/-1.0000,
                   11/-0.5000, 12/-0.4307, 13/-0.3562, 14/-0.3869,
                   15/-0.3333, 16/-0.3155}

\node[strong, anchor=west] at (13.16,1.50) {true $v$};
\fframestems{1.06}{1/1.0000, 2/0.6309, 3/0.5000, 4/0.4307, 5/0.3869,
                   6/0.3562, 7/0.3333, 8/0.3155, 9/-1.0000, 10/-0.6309,
                   11/-0.5000, 12/-0.4307, 13/-0.3869, 14/-0.3562,
                   15/-0.3333, 16/-0.3155}

\node[note, text width=4.70cm] (nA) at (12.60,0.50)
  {$946$ coordinates, $200$ of them non-zero. $\cos(\hat{v},v)$, ranked
   against the other drugs in the same cell line, is $\NIR$.};

\node[note, text width=4.70cm] (nB) at (12.60,-0.60)
  {\DOWN{} has a token position in stage~2 but \emph{no coordinate} here: the
   sign carries the boundary.};

% ================================================== the invariant
\draw[decorate, decoration={brace, amplitude=4pt}, quietink, line width=0.5pt]
  (0.05,-1.80) -- (17.30,-1.80);
\node[anchor=north, font=\footnotesize, text=ink] (nBR) at (8.67,-1.98)
  {and the within-well split-half retrieval reference is \textbf{unchanged}:
   $0.742 \rightarrow 0.743$};

% ---- MEASUREMENT, left in and inert. \fframedebug is undefined in main.tex, so
% this expands to nothing there; figs/_harness_frame.tex \def's it and the
% picture then reports its own extent, the right edge of the target string, and
% the east edge of every node that could plausibly be the widest thing in the
% drawing. All three mattered: the string crossed a divider on the first build,
% and the bounding box ran 1.35pt past the fullwidth measure on a later one.
\ifx\fframedebug\undefined\else
  \path let \p1=(current bounding box.south west),
            \p2=(current bounding box.north east),
            \p3=(c7.east), \p4=(a6.east) in
    \pgfextra{\typeout{FRAME bbox x \x1 .. \x2 ; y \y1 .. \y2}%
              \typeout{FRAME target string ends \x3 ; row (a) ends \x4}};
  \path let \p5=(nW.east), \p6=(nA.east), \p7=(nB.east), \p8=(nBR.east),
            \p9=(nD.east), \p{10}=(nG.east) in
    \pgfextra{\typeout{FRAME east nW \x5 nA \x6 nB \x7 nBR \x8 nD \x9 nG \x{10}}};
\fi
\end{tikzpicture}
\caption[The residual, the target string and the scored vector]{\textbf{The
information was always present; the standard target's tokens barely expose it
--- and the string the model writes is not the vector that scores it.}
\emph{Stage 1.} The drug-specific residual is built by two subtractions. The
control removes the cell's own state, leaving the \emph{shift}; the mean over
the other drugs in the same plate group removes the \emph{generic} programme
that every drug triggers, leaving what makes this drug different. The six
gene profiles are illustrative shapes and carry no measurement --- no axis is
drawn --- but the arithmetic within the figure is exact bar by bar, and the
residual is drawn signed because its sign is what splits the two blocks in
stage~2. \emph{Stage 2.} That same measurement, written two ways. Ranking genes
by expression (a) puts housekeeping genes first, so only $34.6$ of the top $200$
tokens differ between two drugs in the same context; ranking by signed residual
(b) raises that to $118.2$. (Divergence counts quoted at cell-line scope, as in
Table~\ref{tab:divergence}; plate scope is similar, at $120.3$.) Row (b) is the
target string itself: a $100$-gene up block, \DOWN, a $100$-gene down block,
\ENDCELL. \emph{Stage 3.} What $\NIR$ ranks is not that string but a signed
vector over the $946$ panel genes, top-$k$ up positive and top-$k$ down
negative. The stem heights are the exact weights $1/\log_2(\textrm{rank}+1)$, so
the eight drawn per sign are $1.000$, $0.631$, $0.500$, $0.431$, $0.387$,
$0.356$, $0.333$ and $0.316$; which gene sits at which coordinate is
illustrative, and no cosine value is printed, because these two vectors are
drawn rather than measured. Stage~2 and stage~3 are different objects and only
stage~3 is scored. The bracket is the point of the figure --- the within-well
split-half retrieval reference is essentially identical under both encodings, so
no information has been added. What changed is how much of it the token sequence
exposes to the loss --- a property of the encoding, which is not the same as
showing the encoding is what binds (Table~\ref{tab:divergence}).}
\label{fig:frame}
\end{fullwidth}
\end{figure}


% PLACEMENT ONLY. claimbox is `breakable`, and it was breaking here: three lines
% and the accent rule at the foot of one page, the remaining ten at the head of
% the next, so the box read as two fragments and the verdict lost its frame.
% \budgettick reserves 5\baselineskip, which is not enough for a claim this long;
% 16 is the whole box (13 lines plus its padding and two rules).
\needspace{16\baselineskip}%
\begin{claim}
Ranking the target by signed residual exposes $59\%$ of its top-$200$ gene
positions to the drug at cell-line scope, and $60\%$ at the plate scope this
build actually uses, against $17\%$ for the standard target
(\cref{tab:divergence}), at an unchanged within-well split-half retrieval
reference. And the residual is estimable without leaking the evaluation into the
fit: split before fit, generic plate-scoped and drug-meaned, reliability resolved
from the data at the matched null's $95$th percentile ($+0.109$). That is a
statement about the fitting procedure and not about plate structure, which the
scope limit above leaves measured only in expectation. What it establishes is
that the standard target's tokens hide a difference the data contain, and that
this is a fact about ranking genes on absolute expression and not about these
drugs, since $92\%$ of the recoverable between-drug difference returns as soon as
the cell's own control is subtracted. It does not establish that tokenisation is
the binding constraint. The residual package changes five things at once.
\end{claim}

% ...........................................................................
\beatsection{R}{5}{A per-drug constant does not exhaust the target}{sec:res-transfer}
\label{sec:methods-vardecomp}

\begin{question}[5]
The residual says what the model should predict, not how much there is to
predict, and the same training result reads differently against a large reachable
signal and a small one. Two measurements settle the scale: what share of a
perturbation response is drug-specific at all, and how much of that share a table
of per-drug constants would already capture. The first has a clean answer. The
second has an honest one.
\end{question}

% NAMING. The quantity below is the squared Euclidean norm of the shift, and the
% prose says so. shared/inference.py :: energy_shares computes it with nothing
% centred -- denom = s @ s, then (r @ r)/denom, (g @ g)/denom and 2*(r @ g)/denom
% -- so it is a sum of squares about the origin and NOT a variance, whatever that
% function's "Exact variance decomposition" docstring says. Two names in this
% source keep the older signal-processing word on purpose and must not be
% renamed: \label{fig:residual-energy} is a cross-reference target (it prints as
% "Figure 4.12", so the word never reaches the page) and figs/04c-residual-energy
% is an \input target. The reader sees "squared norm" everywhere; only these two
% identifiers still read "energy".
\noindent
The shift's squared norm decomposes exactly into residual $0.621$, generic
$0.390$ and a cross term of $-0.011$, so $62.1\%$ of that squared norm lies
outside the generic programme and $39.0\%$ is the generic programme, the two
essentially orthogonal (\texttt{vardecomp\_matched.json}). The part outside the
generic is the majority of the shift, and the shift is not what a sentence ranks
(\cref{fig:residual-energy}).

% ===========================================================================
% 04c-residual-energy.tex -- the encoding problem in one picture.
% \input by Sections/04c-reencoding.tex. Compiles against thesis_v2/preamble.tex
% and nothing else; no external data, no \includegraphics.
%
% EVERY NUMBER IS SOURCED. Sibling 04c-residual-energy.json carries the key
% paths; the short form:
%   sq.norm  RESULTS_cluster/vardecomp_matched.json -> main.scope
%            residual_energy_share 0.6211205563374411 -> 62.1%
%            generic_energy_share  0.38993489697995193 -> 39.0%
%            cross_energy_share   -0.011055451267510967 -> -0.011
%            energy_shares_sum     1.000000002049882   -> 1.000
%            (the key names say "energy"; the quantity is the squared L2 norm)
%            (text, all three agreeing; line numbers drift, so each carries the
%             phrase to search for:
%             04c-reencoding.tex:301-303  "decomposes exactly into residual"
%             01-Introduction.tex:313-315 "Decomposing the treated-minus-control"
%             05-Limitations.tex:240-242  \paragraph{Frame coverage}
%             -- verify by the phrase, not the number, before quoting a line)
%   tokens   RESULTS_cluster/target_divergence.json, K=200
%            scopes.cellline.full.divergence_mean     34.594876910336524 -> 34.6 (17%)
%            scopes.cellline.shift.divergence_mean   111.4231960024091   -> 111.4 (56%)
%            scopes.cellline.residual.divergence_mean 118.24622167432057 -> 118.2 (59%)
%            (text: tab:divergence, 04c-reencoding.tex:42-44, the three body
%             rows of that table -- agrees)
%   the caption's 120.3 is scopes.plate.residual.divergence_mean 120.29533990815328
%   and is quoted in prose only; it is NOT drawn, because a plate-scope bar in a
%   cell-line-scope block is exactly the cross-block subtraction this layout is
%   built to prevent.
%
% LAYOUT NOTES, all forced by a build:
%   * two SUB-BLOCKS, one axis, one divider. The blocks are DIFFERENT QUANTITIES
%     (a squared-norm share and a count of token positions); the divider, the
%     1.68cm gap, the 2:1 bar heights (4.8mm vs 2.4mm) and the unit sentence
%     are what stop a reader subtracting across them. Nothing spans the divider.
%   * TOTAL WIDTH 14.191cm = 403.775pt against \textwidth = 404.02913pt (142.0mm,
%     measured in _v.log; preamble.tex:59 records the 132->142mm revision). The
%     figure fills 99.94% of the measure and needs no fullwidth, so the page it
%     lands on may still carry a margin note. It was previously laid out for the
%     132mm block, drew 348.3pt = 122.4mm and under-filled by 19.6mm.
%     Geometry: \X 4.900cm label gutter + \L 9.190cm per unit share; the right
%     end of the overhang is \X+1.011*\L = 14.19109cm. Picture box measured by
%     _scratch/figs/_harness_resenergy_w.tex, which reports
%     PICW = 403.77531pt PICH = 124.33833pt in _scratch/_harness_resenergy_w.log.
%   * ROW LABELS ARE ONE LINE EACH. Measured at \scriptsize in Pagella:
%     "residual --- the drug-specific target" 122.69pt = 4.312cm is now the
%     widest ("the perturbation response" measures 91.40pt), so the label column
%     stays at 4.65cm. Two-line labels at 0.62cm row pitch collide, which is why
%     the width is measured and not guessed.
%   * \bx must expand WITHOUT braces: it is used bare, as (\bx{0.5},y), and
%     inside an expression, as ({\bx{0.621}+0.14},y). A braced body makes the
%     second form "{{...}+0.14}" and pgfmath aborts with "Missing number".
%   * the 62.1% and 39.0% labels sit INSIDE/AFTER their segments; the token-block
%     labels sit AFTER the accent segment, because the 17% segment is 1.56cm
%     wide and will not hold its own label.
%   * THE HALF-WAY RULE IS DRAWN ONLY WITHIN THE TWO BAR BANDS. The previous
%     version ran it floor-to-ceiling and struck through both grey annotations
%     ("shift's" in the cross-term line, "positions" in the unit line); the old
%     header's guarantee that the rule "passes through none of them" guarded the
%     bold in-bar labels only and never covered the prose. Measured clearances
%     at bx{0.5} = 9.495cm, all in the bar bands: "62.1% drug-specific" ends at
%     7.50cm (1.99cm clear), "34.6 (17%)" ends at 8.04cm (1.46cm clear), and
%     "111.4 (56%)", "118.2 (59%)" and "generic 39.0%" all begin right of the
%     rule at 10.19, 10.46 and 10.75cm. The rule is drawn last, over the bars.
%     The two segments are given the same vertical extent as the upright ink
%     rule at 1.000, so the figure's two verticals rhyme instead of competing.
%   * THE "half" LABEL SITS UNDER THE TOKEN BLOCK, not at the top of the figure.
%     The half-way mark does its work below the divider, where 17% / 56% / 59%
%     are read against it; at the top it was as far from that reading as the
%     picture allowed.
%   * the generic segment runs on from where the residual segment ends, out to
%     1.011, with the ink rule at the shift's own squared norm: the
%     decomposition is exact, and drawing the -0.011 overhang rather than
%     describing it keeps the two shares from being silently renormalised to
%     sum to one.
%   * the cross-term annotation is anchored EAST and grows leftwards. Anchored
%     WEST at the right end of the bar it walked 51.1pt off the page.
% ===========================================================================
\begin{figure}[htbp]
\centering
\begin{tikzpicture}[x=1cm, y=1cm, font=\small]
  \def\X{4.900}                  % left edge of every bar
  \def\L{9.190}                  % bar length for 100%
  \newcommand{\bx}[1]{\X+#1*\L}  % value -> x   (no braces: see header)
  % PIN THE BOUNDING BOX. Left free, the picture measures 409.59pt and runs
  % 5.56pt over the block: the rowname nodes are anchor=east at 4.75 with a
  % 4.65cm text width, so their inner sep hangs ~0.09cm left of x=0, and the
  % "half" descender hangs below. Neither carries ink -- the row labels are
  % flush right and set single-line, so their ink starts at x=0.07cm. Pinning
  % to [0, \bx{1.011}] therefore clips nothing and makes the drawn width equal
  % the stated 14.191cm = 403.775pt exactly.
  \useasboundingbox (0,-0.25) rectangle ({\bx{1.011}},4.12);
  \tikzset{
    rowname/.style={anchor=east, text width=4.65cm, align=flush right,
                    font=\scriptsize, text=quietink},
    inlab/.style={anchor=west, font=\scriptsize\bfseries, text=white},
    outlab/.style={anchor=west, font=\scriptsize\bfseries, text=accent},
    greylab/.style={anchor=west, font=\scriptsize, text=ink}}

  % ========== SUB-BLOCK 1: the response, by squared norm. Bars 4.8mm. ======
  \node[rowname] at (4.75,3.70) {the perturbation response};
  \fill[accent]      (\X,3.46)          rectangle (\bx{0.621},3.94);
  \fill[rulegrey!45] (\bx{0.621},3.46)  rectangle (\bx{1.011},3.94);
  \node[inlab]   at ({\X+0.14},3.70)          {$62.1\%$ drug-specific};
  \node[greylab] at ({\bx{0.621}+0.14},3.70)  {generic $39.0\%$};

  % the shift's own squared norm, and the cross term as the overhang past it
  \draw[ink, line width=0.7pt] (\bx{1.0},3.32) -- (\bx{1.0},4.08);
  \node[anchor=north east, text width=7.6cm, align=flush right,
        font=\scriptsize, text=quietink] at ({\bx{1.011}},3.28)
    {the upright rule marks the shift's own squared norm; the overhang past it
     is the cross term $-0.011$};

  % --- the divider: the two sub-blocks are different quantities ------------
  \draw[rulegrey, line width=0.4pt] (0,2.44) -- ({\bx{1.011}},2.44);
  \node[anchor=north west, font=\scriptsize\itshape, text=quietink] at (0,2.38)
    {above: shares of the shift's squared norm \quad below: shares of $200$
     token positions};

  % ========== SUB-BLOCK 2: the target, by token position. Bars 2.4mm. ======
  % r1 -- the standard target
  \fill[rulegrey!25] (\X,1.54) rectangle (\bx{1.0},1.78);
  \fill[accent]      (\X,1.54) rectangle (\bx{0.17},1.78);
  \node[rowname] at (4.75,1.66) {full profile --- the standard target};
  \node[outlab]  at ({\bx{0.17}+0.14},1.66) {$34.6$ \;($17\%$)};

  % r2 -- the shift
  \fill[rulegrey!25] (\X,0.92) rectangle (\bx{1.0},1.16);
  \fill[accent]      (\X,0.92) rectangle (\bx{0.56},1.16);
  \node[rowname] at (4.75,1.04) {shift (treated $-$ control)};
  \node[outlab]  at ({\bx{0.56}+0.14},1.04) {$111.4$ \;($56\%$)};

  % r3 -- the residual
  \fill[rulegrey!25] (\X,0.30) rectangle (\bx{1.0},0.54);
  \fill[accent]      (\X,0.30) rectangle (\bx{0.59},0.54);
  \node[rowname] at (4.75,0.42) {residual --- the drug-specific target};
  \node[outlab]  at ({\bx{0.59}+0.14},0.42) {$118.2$ \;($59\%$)};

  % --- the half-way rule, drawn LAST so the bars do not bury it, and drawn
  % ONLY within the two bar bands so it crosses no prose. See header.
  \draw[ink, line width=0.5pt, densely dotted] (\bx{0.5},3.32) -- (\bx{0.5},4.08);
  \draw[ink, line width=0.5pt, densely dotted] (\bx{0.5},0.18) -- (\bx{0.5},1.90);
  \node[anchor=north, font=\scriptsize\itshape, text=ink] at (\bx{0.5},0.14) {half};
\end{tikzpicture}
\caption[The drug is most of the response, little of the
target]{\textbf{The drug is the majority of the perturbation response and almost
none of the standard target's tokens.} Above the rule: the shift's squared
norm decomposes exactly into a drug-specific residual, a generic programme and
a cross term, so $62.1\%$ of what a perturbation does to a cell is specific to
the drug. Below it: of the $200$ positions a cell sentence spends on the
standard target, only $34.6$ carry a different gene for two drugs in the same
context --- the other $83\%$ are the same string. Ranking the same measurement
by signed residual raises that to $118.2$, and at the plate scope this build
trains in the residual row reads $120.3$ ($60\%$). The two halves of the figure
are different quantities, a squared-norm share and a count of token positions;
they share the axis and the half-way rule and nothing else, and no difference
between them is a number. What the picture states is an ordering: the signal is
a majority of the response, the standard encoding exposes a small minority of
it to the loss, and the residual encoding puts the exposed fraction back on the
majority side (\cref{tab:divergence}).}
\label{fig:residual-energy}
\end{figure}


That share has not previously been computed for this task, and its frame is not
the trained-on one. The stored generic averages over conditions and keeps the
condition's own contribution, where the training frame averages over drugs and
leaves the condition's own drug out. Self-inclusion scales a residual by
$(m-1)/m$ in a plate group of $m$ conditions, so that frame difference depresses
the share quoted here; it does not inflate it. The weighting difference has no
established direction. A second bias runs the other way. A condition's pseudobulk
sampling error enters its own shift whole and is averaged down in the group mean,
so it lands in the residual term and inflates the estimated share at this cell
budget (\cref{tab:scale}, and the reliability line above). The two oppose each
other, the net is not established, and no noise-corrected share is estimated
here.

\paragraph{Two components ask for different work} Writing
\begin{equation}
 \resid(d,c) \;=\; \beta(d) \;+\; \kappa(d,c) \;+\; \epsilon,
 \qquad
 \Tcoef \;=\; \frac{\operatorname{var}\beta}{\operatorname{var}\beta + \operatorname{var}\kappa},
\end{equation}
the per-drug constant $\beta(d)$, the drug \emph{main effect}, is what a lookup
table reproduces with no model at all.\seealso{\cref{sec:res-lookup} builds that
lookup and scores it.} The remainder $\kappa(d,c)$ is the only component a
conditional model could win on. This is DrEval's naive predictor~\cite{dreval} in
transcriptomic space. Their \texttt{NaiveMeanEffectsPredictor} predicts
$\hat{y}_{ij} = \mu^c_i + \mu^d_j - \mu$, cell-line and drug main effects with no
interaction term, and $\textrm{generic}(c) + \beta(d)$ is that quantity lifted
from a scalar potency to a $946$-dimensional profile.

Two warnings attach before any number is quoted. $\kappa$ is written as a
drug\,$\times$\,cell-line interaction because the decomposition assumes one, but
the estimator used does not identify such a term, so $\kappa$ is called the
\emph{context-specific remainder} and $\Tcoef$ a disattenuated cross-context
residual-similarity coefficient, not a variance share. And $1 - \Tcoef$ does not
measure the headroom above the additive baseline. The cross-line pairs entering
it may also differ in dose and in treatment well, and each residual is taken
against a generic its own drug helped define, so it bounds that headroom from
above at best.

% PLACEMENT: this box carries a display equation and cannot break usefully. At
% the previous setting it opened at the foot of the page and tcolorbox broke it
% immediately after the title bar, stranding "DEFINITION The transfer
% coefficient" alone at the page bottom with every line of the body overleaf.
% \budgettick's own \needspace{5\baselineskip} is not enough to catch that; the
% box is ~13 lines tall, so ask for that and let it move whole.
\needspace{12\baselineskip}
\begin{defn}{The transfer coefficient}
With $\num{286}$ of $\num{287}$ treatments in exactly one well there is
no replicate from which to separate $\operatorname{var}\kappa$ from
$\operatorname{var}\epsilon$, so the variance-components model is unidentified on
its own. The cosine route buys identification through split-half reliability:
\begin{equation}
 \hat{\Tcoef} \;=\; \frac{\cos\!\big(\resid(d,c_1),\, \resid(d,c_2)\big)}
 {\sqrt{\rho(d,c_1)\,\rho(d,c_2)}},
 \qquad
 \rho(d,c) \;=\; \frac{2\cos(\resid_A, \resid_B)}{1 + \cos(\resid_A, \resid_B)},
\end{equation}
where $\resid_A$ and $\resid_B$ are disjoint halves of condition $(d,c)$ and
$\rho(d,c)$ is the Spearman--Brown reliability of that condition's residual at
its full cell count.
\end{defn}

Dividing out the reliability is essential, since measurement noise alone
depresses the raw cross-line cosine and would otherwise be read as interaction.
Two cautions ride with it. Spearman--Brown is derived for correlations between
parallel halves, and what licenses it on a cosine here is the planted-truth
validation in \cref{sec:appendix}. And the reliability entering it is a
within-well split-half precision, optimistic relative to true biological
reproducibility. That optimism sits in the denominator, where an over-large
$\rho$ makes $\hat{\Tcoef}$ too small, so $1 - \hat{\Tcoef}$ is, if anything,
\emph{overstated}. An earlier draft had that direction the wrong way round.

\paragraph{The control that had to be replaced} The first run compared $\Tcoef$
against a negative control of different drugs on the same plate, which returned
$+0.485$ against $\Tcoef = 0.511$, and was correctly declared
void.\corrmark{the same-plate negative control is voided, not matched; the
structure-matched control replaces it} That control differs from the estimand in
two ways at once, different drug \emph{and} same cell line, and plate structure
surviving a cell-line-scoped generic inflates same-plate comparisons while
leaving cross-line comparisons untouched. The $+0.485$ is voided, not matched
against the separate $-0.017$ plate-scoped diagnostic. Only a structure-matched
control, breaking the drug match alone, bounds the estimand.

% ...........................................................................
\begin{table}[htbp]
\begin{threeparttable}
\caption{\textbf{The five checks accompanying the transfer-coefficient
estimate, and what each one bounds.} At $\kappa = 0$ the \emph{uncorrected}
cross-line cosine is only $0.734$, which read naively reports $27\%$ interaction
where the truth is zero, while the disattenuated estimator recovers the truth to
within $0.001$. That recovery licenses the reference column of
\cref{tab:transfer}.}
\label{tab:tcontrols}
\small
\begin{tabularx}{\linewidth}{@{}>{\raggedright\arraybackslash}p{0.34\linewidth}X@{}}
\toprule
\thead{Check} & \thead{What it bounds} \\
\midrule
disjoint control halves & shared-control inflation: $+0.242$ at cell-line scope,
  $-0.000$ at plate \\
structure-matched negative control & the estimand itself; returns $-0.008$ \\
simulated $\kappa{=}0$ null & estimator bias: returns $1.001$ where the truth is
  $1.000$ \\
dose strata on molar doses & a reference scale: the same drug and line at
  another dose \\
leave-one-drug-out generic & train/held-out target asymmetry
  (\cref{sec:methods-residual}) \\
\bottomrule
\end{tabularx}
\begin{tablenotes}[flushleft]\footnotesize
\item An earlier cache carried the well identifier in the dose field, so what was
labelled dose transfer compared physical wells, not concentrations.
\end{tablenotes}
\end{threeparttable}
\end{table}

% ...........................................................................
\begin{table}[htbp]
\begin{threeparttable}
\caption{\textbf{The transfer coefficient, estimated on real molar doses with
dyadic cluster-robust intervals.} The matched negative control sits on zero. The
same-plate control that produced the earlier false alarm reads $+0.485$ under a
cell-line-scoped generic. $\Tcoef$ is a cross-context similarity coefficient, not
a variance share, and $1 - \Tcoef$ must not be read as a
drug\,$\times$\,cell-line interaction share.}
\label{tab:transfer}
\small
\begin{tabularx}{\linewidth}{@{}%
  >{\raggedright\arraybackslash}p{0.28\linewidth}
  >{\raggedright\arraybackslash}p{0.13\linewidth}
  >{\raggedright\arraybackslash}p{0.20\linewidth}
  >{\raggedright\arraybackslash}X@{}}
\toprule
\thead{Quantity} & \thead{Estimate} & \thead{$95\%$ interval} & \thead{Reference} \\
\midrule
$\Tcoef$ (cross-context, plate-scoped generic) & $0.553$ & $[0.514, 0.593]$
  & simulated $\kappa{=}0$ null $= 1.001$ \\
similarity scale NOT shared across contexts, $1-\Tcoef$ & $44.7\%$
  & $[40.7\%, 48.6\%]$ & mixes cell line, dose, well and estimator \\
$\Tcoef$ across dose (same drug, same line) & $0.707$ & $[0.632, 0.782]$
  & reference scale \\
structure-matched negative control & $-0.008$ & spanning zero
  & must be $\approx 0$ \\
\bottomrule
\end{tabularx}
\end{threeparttable}
\end{table}

\paragraph{What the coefficient is, precisely} $\Tcoef$ is a descriptive
sensitivity analysis and not a load-bearing estimate, and the reason sits in its
configuration. The scoped calculation retains conditions at a reproducibility
threshold of $0.2$, and it carries every caution above at once: a
\emph{self-including} generic, pairs free to differ in \emph{dose} and in
treatment well, and an optimistic reliability in the denominator. That last one
pushes $\hat{\Tcoef}$ down, so $1 - \Tcoef$ errs high and bounds from above what
a per-drug constant leaves on the table, without measuring it.

\paragraph{Dose is the milder perturbation} Holding drug and cell line fixed and
changing only the molar concentration, transfer is $0.707$ $[0.632, 0.782]$ on
$\num{363}$ pairs against $0.553$ $[0.514, 0.593]$ pooled cross-context; the
intervals do not overlap. That contrast is between one clean axis and a mixture,
since the cross-context arm differs in well and dose as well as cell line, but it
is the closest this section comes to a single-axis comparison.

The same ordering served as the weaker of the two scope diagnostics. At cell-line
scope it inverts. The dose arm falls to $0.4846$, a different quantity from the
voided same-plate control, which coincidentally reads $0.4853$, while the
cross-line coefficient holds at $0.511$. That is the wrong way round for dose
against cell line, but the expectation is a prior, so it is used as a
plausibility check. The second diagnostic needs no prior. At plate scope the
shared-control and split-control builds of $\Tcoef$ agree to $-0.000$, while at
cell-line scope they differ by $+0.242$, plate structure leaking through a
cell-line-scoped generic. That is what the scope decision rests on, and it is why
the cell-line-scoped build cannot corroborate when it returns a similar
$\Tcoef = 0.511$; that build is the one declared void, its own same-plate
negative control reading $+0.485$.

\paragraph{Whether anything could learn the remainder} If a cell line modulated
drug responses in a consistent direction, that direction would be estimable from
the line's own cells. Writing $\kappa(d,c) = \resid(d,c) - \hat{\beta}(d)$, the
test compares $\cos(\kappa(d,c), \kappa(d',c))$ for different drugs in the same
cell line against the same quantity across different cell lines, with
$\hat{\beta}$ estimated leave-one-cell-line-out; estimated over all lines it
would contain $\resid(d,c)$ itself, and subtracting a mean containing its own
term induces a $-1/(m-1)$ correlation that biases the signal arm downward by
construction. Both arms are disattenuated by $\kappa$'s own split-half
reliability.

Averaged over lines, no consistent per-line direction shows up. Under the
plate-scoped generic the same-cell-line arm reads \ci{-0.0149}{-0.0251}{-0.0050}
and the different-cell-line arm \ci{+0.0003}{-0.0071}{+0.0081}, an excess within
line of $-0.0153$ for which no interval was computed. The intervals come from a
cell-line-clustered bootstrap over $\num{38}$ and $\num{42}$ clusters, on
$\num{4000}$ same-line and $\num{3884}$ cross-line pairs over $\num{1282}$
conditions. Neither arm shows the positive within-line excess a per-line
direction would require, and the negative sign of the same-line arm is not read
as evidence: same-line pairs that also share a plate share a generic, and that
common subtraction pushes their cosines negative by construction, where
cross-line pairs share no generic at all. It is the construction the
different-drug same-plate diagnostic reads at $-0.017$ above. A
cell-line-scoped counterpart was not computed in the canonical run, and the row
an earlier draft carried for it is withdrawn.

Two negative results do not add up to irreducibility: this one, and a direct
test of whether mechanism-matched drugs deviate from their per-drug constants in
the same way within a cell line, which returned nothing detectable. The second
bounds its own power severely. $\kappa$'s split-half reliability in this atlas is
$0.195$ at the $0.05$ reliability floor ($0.324$ at the stricter $0.20$ floor),
and same-mechanism drugs are so heavily co-plated that only $3$--$7\%$ of
matched pairs survive a different-plate requirement ($1\%$ at the stricter
floor). Resolving $\kappa$'s structure would take more cells per condition or a
plating design that crosses mechanism with plate.

\begin{claim}
There is a conditional target that a per-drug constant cannot express. Two tests
looked for learnable structure in it and found none, and neither settles whether
there is any. $\kappa$'s split-half reliability is $0.195$ at the floor the
mechanism test runs at, and only $3$--$7\%$ of mechanism-matched pairs survive a
different-plate requirement, so its structure is unresolved here rather than
\shownzero. Its size is quoted as what it is: $\Tcoef = 0.553$
\nolinebreak[3]$[0.514, 0.593]$ is a disattenuated cross-context similarity
coefficient, so $44.7\%$ $[40.7\%, 48.6\%]$ of \emph{that similarity scale}
remains unshared across the compared contexts. It is not a variance share and not
a drug\,$\times$\,cell-line interaction, since the estimator pools pairs
differing in cell line, dose, well and target estimator, so it bounds what a
per-drug constant leaves on the table without establishing that a model could
reach any of it.
\end{claim}

% ...........................................................................
\beatsection{R}{5}{Re-encoding produces prompt sensitivity}{sec:res-encoding}

\begin{question}[5]
Every target change tried so far left the model indifferent to the drug named in
its prompt. Does training on the residual change that? The instrument is the
stratified scramble, the same prompt with another drug's name and mechanism
substituted, read as a ladder of strata because the swapped-in drug is itself a
drug the model knows and so is not a neutral comparison. The model moves. It is
not the only target change that makes it move.
\end{question}

\noindent
On the $\num{300}$ trained conditions of the residual-frame evaluation the model
beats the geometry-defined \texttt{orth} scrambled partner by
\ci{+0.0898}{+0.0291}{+0.1506}, clustered on the drug and the treatment well
(\texttt{re\_v3.json}).\gloss{scramble}{The same prompt with the drug's name and
mechanism replaced by another drug's, in three strata by residual similarity.}
That is the first non-null drug effect in this project, first in the order the
experiments were run, and not the only target construction that produces one.

% ...........................................................................
\paragraph{Why the comparison set is not exchangeable, and what follows} The
comparator is not a null. Its own score is a monotone function of which drug it
was handed, which makes it an active treatment; \cref{sec:res-generalisation}
measures the size of that. Two consequences fix the reading protocol here. The
series is quoted as a \emph{gradient}, since a metric artefact would not order
itself by swap dissimilarity; and where a single magnitude is needed it is the
\texttt{orth} one, the least contaminated comparison available and not an
independently established null, that stratum having been defined using true
residual geometry.

\paragraph{Why the strata are cut by response geometry and not by mechanism} A
scramble that swaps in a drug with a similar true response is a weak test, and
the annotation that looks like it should prevent that does not. In the
characterisation behind \cref{tab:scale}, two drugs sharing a \texttt{moa-fine}
annotation sit a mean distance $2.322$ apart against $2.377$ for two drugs that
do not, $\num{1244}$ same-mechanism against $\num{17563}$ different-mechanism
pairs, a ratio of $0.977$. Same-mechanism drugs in Tahoe are just over two per
cent closer to each other than arbitrary pairs are, so ``different mechanism''
does not imply ``different response''. The strata are therefore cut by measured
residual cosine within the cell line. (Mechanism does carry signal in the
residual frame, \cref{sec:res-scope}, a different frame from the full profiles
measured here.)

\paragraph{The gradient, on the earlier target build} The stratified scramble of
\cref{sec:methods-controls} gave a gap that grew with swap dissimilarity
(\cref{fig:repair}), \ci{+0.0351}{+0.0061}{+0.0656},
\ci{+0.0716}{+0.0371}{+0.1059} and \ci{+0.1429}{+0.1112}{+0.1787}
across \texttt{near}, \texttt{orth} and
\texttt{opposite}, each differenced against its own stratum-specific partner.
These are the residual arm's $1400$-token run; the two control arms beside them
were scored at the $600$-token cap of \cref{sec:methods-data}, where their
stratum gaps run from $-0.021$ to $+0.026$ and were read at the time as flat and
null. That reading holds for the single-cell arm at both budgets. It does not
hold for the optimal-transport arm, whose $600$-token gaps are themselves
monotone in swap dissimilarity.

\begin{figure}[htbp]
\centering
\includegraphics{fig-repair}
\caption[The gap grows with swap dissimilarity]{\textbf{On the earlier target
build, the residual model's advantage over its own scramble grew with how
dissimilar the swapped-in drug's residual was; the single-cell control scored
through the identical pipeline stayed flat and null, while the optimal-transport
control did not.} Residual frame throughout. The horizontal axis is measured
swap dissimilarity, $1-\cos$ between the real and the swapped-in drug's residual,
byte-identical across all four runs, so the three stratum positions are exact;
the connecting lines are guides between three measured strata and not a sampled
continuum, and the four series are nudged sideways so that their intervals do not
overlap. The three solid series are at a matched $600$-token
budget; the residual model's $1400$-token run is the dotted overlay, since the
budget alone roughly doubles the gap on the two strata where it is positive at
$600$ ($+0.0349 \rightarrow +0.0716$ on \texttt{orth},
$+0.0716 \rightarrow +0.1429$ on \texttt{opposite}). Intervals are $95\%$
bootstrap clustered over cell lines, $n = 250$ conditions per point. The
\texttt{near} stratum is the weakest test by construction, since where the
swapped-in drug's true residual is similar, an unchanged prediction is correct.}
\label{fig:repair}
\end{figure}

\paragraph{Three diagnostics that need no scramble} All three agreed on that
build, and the residual model led on all three. The correlation between a
condition's reproducibility and model $\NIR$ was $+0.144$, the model doing better
where the drug effect is real, against $-0.117$ and $-0.216$ for the controls.
Prediction diversity was highest for the residual model. And
$\cos(\textrm{prediction}, \textrm{own truth})$ exceeded
$\cos(\textrm{prediction}, \textrm{other drugs})$ for all three arms, the
residual model's margin, $+0.084$, being roughly twice the optimal-transport
arm's $+0.049$ and three times the single-cell arm's $+0.026$.

\paragraph{One measurement lesson before the withdrawal} The $600$-token cap
roughly halved the effect at every stratum where it was positive there, so
$+0.0716$ appears twice in this section, once as the $600$-token
\texttt{opposite} gap and once as the $1400$-token \texttt{orth} gap, and the two
are not the same quantity. Re-run at $1400$ tokens, the single-cell control stays
flat, $-0.0128$ \nolinebreak[3]$[-0.0458, +0.0202]$ against its own
\texttt{orth} partner. The optimal-transport control does not.

% PLACEMENT ONLY. Worst case of the breakable box: nine lines of the retraction
% at the foot of one page, then \cref{fig:repair} taking the whole top of the
% next as a deferred float, then the remaining six lines under it -- the box read
% as two boxes with a figure between them. 17 lines is the box entire, so it
% either sets whole here or starts the next page whole, and the figure can no
% longer land inside it.
\needspace{17\baselineskip}%
\begin{withdrawn}{``OT achieved nothing'', \cref{sec:res-epsilon}}
At $600$ tokens the optimal-transport arm's \texttt{orth} gap was
\ci{+0.0144}{-0.0051}{+0.0338} and spanned zero; at $1400$ it is
\ci{+0.0292}{+0.0144}{+0.0441} and does not. The token budget was suppressing a
real effect in that arm, the same lesson the residual arm taught, applied to a
control previously read as inert. The $600$-token evidence pointed the same way
and was read too cautiously; that arm's stratum gaps there were $+0.0073$,
$+0.0144$ and $+0.0263$, monotone in swap dissimilarity like the residual arm's,
with only \texttt{opposite} excluding zero. The consequence reaches backwards.
The optimal-transport target is one point in the three-point comparison of
\cref{sec:res-epsilon}, and that negative holds only at the budget and comparator
it was measured with. At a matched budget and against the geometry-defined
\texttt{orth} partner, whose true residual is on average orthogonal to the
target's ($\cos = +0.035$), the optimal-transport point is not null. The
single-cell point is the one whose interval still spans zero. The geometry-based
selection does not independently establish a null in either case.
\end{withdrawn}

\paragraph{The comparison that costs this section its strongest phrasing} The two
arms' pooled gaps against their own \texttt{orth} partners overlap almost
entirely (\texttt{ctrl1400\_summary.json}), so re-encoding the target produces
prompt sensitivity but is \emph{not} the only target change that does.

Three limits stop that comparison being read as an ordering. The arms are scored
in different frames, so only each arm's gap against its own comparator is
comparable; they are separate evaluations over different condition sets, so this
is an arm-level and not a paired comparison; and those sets differ in split
composition by more than the arms differ from each other. The residual figure
pools the three splits of \cref{tab:generalisation}, whose own gaps are
$+0.0898$, $+0.0332$ and $-0.0315$, while the optimal-transport evaluation
carries no split labels at all. A spread of that size across splits, against a
$0.0069$ difference between the two pooled numbers, means the near-equality does
not order the arms in either direction.

\paragraph{The slope carries its own null} Regressing each prediction's $\NIR$
on the cosine between the named drug's truth and the target's truth gives a
slope whose null is zero for a model that ignores the drug it is told. Its
points are the three scrambled partner prompts and nothing else, the target-drug
prompt never among them, and \cref{sec:res-generalisation} fits it on all three
splits. The fit is pooled least squares over every point in a split with a
single intercept, three points per condition, no condition-level intercepts
absorbed, so the zero null is exact within a condition while the pooled slope
also carries a between-condition component. What it closes is the freedom to
pick a comparator, since it spends all three strata rather than one.

\begin{claim}
Re-encoding the target produced the project's first non-null drug effect. On
trained conditions the model beats the geometry-defined \texttt{orth} partner by
\ci{+0.0898}{+0.0291}{+0.1506} and beats chance outright, with no comparator arm
at all, by \ci{+0.0981}{+0.0407}{+0.1554}, both clustered on the drug and the
treatment well, the drug being the axis the claim rides on
(\cref{sec:res-generalisation}). The comparator-free slope of retrieval against
the named drug's truth is \ci{+0.3728}{+0.2797}{+0.4660}, and zero is its null for a
drug-blind model, so no choice of comparator manufactures it, though it is fitted over the
scrambled partner prompts alone, with the target-drug prompt never a point on the
line. It is first in sequence and not uniquely responsible: at a matched
$1400$-token budget the optimal-transport target clears zero against its own
\texttt{orth} partner too (\ci{+0.0292}{+0.0144}{+0.0441} over $\num{248}$
conditions that carry no split labels, against the residual arm's
\ci{+0.0361}{+0.0159}{+0.0564} pooled over $\num{1394}$ that mix trained and
held-out splits), on a separate evaluation in a different frame, so the two are
not ordered against each other. The model has been shown to respond to the drug;
whether that response is any use is the next question.
\end{claim}

% ...........................................................................
\beatsection{R}{6}{Responding to the drug is not predicting it}{sec:res-generalisation}
\label{sec:methods-holdout}

\begin{question}[6]
A model that moves when the drug token moves has cleared a low bar, because
reproducing a training average requires that much too. Everything scored so far
came from conditions the model trained on, so what has been shown is that the
token is consulted, not that anything transferable was learned. Separating the
two needs a holdout, and the
atlas's own randomisation decides which one means anything: whole treatment
wells. Prompt sensitivity survives that test. Whether the residual itself
transfers comes back unsettled.
\end{question}

\noindent
The physical layout of \cref{sec:methods-data} forces the choice. A (drug, cell
line) holdout, the field's standard Leave-Pairs-Out split and the one an earlier
version of this work used, leaves the held-out condition's treated well sitting
in the training set through its other $\approx 48$ cell lines. Measured against
that holdout, $\num{124}$ of the $\num{255}$ treated wells for which split
information exists placed a held-out condition and a training condition in the
same well.\seealso{\cref{sec:res-lookup}: what a held-out gap has to
beat.} That estimand measures within-atlas matrix completion, not independent
generalisation, and no processing of the targets changes it.

Holding out whole wells is the only assignment-respecting choice the design
offers, and it answers a narrower question, generalisation to an unseen treatment
well of a seen drug. Fixed-dose cross-context transfer, present in one line and
absent in another, is structurally unidentifiable in this atlas.

\paragraph{What each split holds} \texttt{train} takes every condition of the
remaining wells, $\num{4999}$ before the reliability line and
\needsaudit[tally against the gate's retained count]{2{,}523} after it.
\texttt{unseen\_combo} holds $\num{36}$ treated wells out whole ($\num{904}$
conditions, $\num{900}$ after inventory filters), and it never takes a drug's
last training well; without that guard a two-well drug could silently become an
unseen drug while sitting in the unseen-combination arm, and the two arms would
measure the same thing. \texttt{unseen\_drug} holds every condition of
$\num{10}$ held-out drugs ($\num{714}$ conditions, $\num{711}$ after inventory
filters). Zero wells contribute conditions to both sides.

\texttt{unseen\_drug} is a control and not a third result. Tahoe fine-tuning
never sees the token, so a large effect there would indict the comparator, not
report a biological finding. ``Unseen'' means held out from fine-tuning with
mechanism supplied, not unknown to the pipeline; the Pythia-1b backbone was
pretrained on natural language and the prompt still carries the
\texttt{Mechanism:} field.

Two warnings ride with the design. DrEval reports, citing~\cite{partin}, that
apparently good performance on pair-holdout splits can be reached simply by
predicting each drug's average response across the training cell lines, which is
why the baseline ladder of \cref{sec:res-lookup} is not optional. And the
evaluation draws under a per-split quota, because a uniform draw reproduces the
pool's proportions and starves the split the experiment exists to test; in a
first run,
$\num{250}$ constructed held-out combinations yielded only $\num{33}$ scored
conditions. Under the quota the evaluation scores $\num{1394}$ conditions in all,
$\num{300}$ trained and $\num{199}$ unseen-drug, the rest unseen combinations,
almost the whole pool of $\num{900}$ that arm affords, while
\texttt{unseen\_drug} remains a quota of the $\num{711}$ available.

% ...........................................................................
\paragraph{The first result of this design is about the comparator} It is a
finding and not an erratum, and it reframes every gap quoted so far. Across the
three swap strata the scrambled model's $\NIR$ is $0.550$, $0.501$ and $0.423$ as
the named partner's true residual moves from aligned ($\cos = +0.22$) through
orthogonal ($0.00$) to anti-aligned ($-0.20$) with the target. The
\texttt{near} and \texttt{opposite} scores differ from chance in opposite
directions; the geometry-defined \texttt{orth} partner scores $0.501$ and is
empirically near chance on these conditions. That makes \texttt{orth} the least
contaminated available comparator, not an independently established null, since
its stratum was defined using true residual geometry.

The consequence is structural. A gap measured against such an arm decomposes as
\begin{equation}
 \gap \;=\; \underbrace{\big(\NIR_{\textrm{model}} - 0.5\big)}_{\textrm{how much naming the truth helps}}
 \;+\; \underbrace{\big(0.5 - \NIR_{\textrm{scramble}}\big)}_{\textrm{how much naming the lie hurts}},
\end{equation}
and only the first term is wanted. The second is not a rounding error. Pooled
over all $\num{1394}$ scored conditions the model's own advantage over chance is
$+0.0367$ and the \texttt{opposite} comparator's deficit is $+0.0772$, so $68\%$
of the $+0.1139$ gap that arm reports is the comparator failing and not the model
knowing. Against \texttt{orth} the second term is approximately zero here, so
every gap in \cref{tab:generalisation} is measured against that partner, its
magnitude conditional on the observed near-chance score.

The diagnostic is free and not specific to this experiment. Score the comparator
arm against chance; if its score moves with which label it was handed, the gap it
defines is not a measurement of the model. A swap drawn to be maximally
dissimilar, the obvious way to build a strong negative control, is exactly the
swap that maximises the second term.

% ...........................................................................
\begin{table}[htbp]
\begin{fullwidth}
\begin{threeparttable}
\caption{\textbf{Whole-well holdout on the corrected evaluation, with the truth
verified against the build's own fit digest and the full held-out population
scored.} The two right-hand columns carry the SAME point estimate and differ only
in what they treat as the unit of generalisation. On \texttt{unseen\_combo} the
gap excludes zero clustered on cell line and does not clustered on drug, which is
why that row is reported as \notestablished. The \texttt{unseen\_drug} row
corroborates but does not prove: its \texttt{orth}-partner gap is negative while
its \texttt{opposite} gap is $+0.0379$.}
\label{tab:generalisation}
\small
\begin{tabularx}{\linewidth}{@{}%
  >{\raggedright\arraybackslash}p{0.16\linewidth}
  >{\raggedright\arraybackslash}p{0.18\linewidth}
  >{\raggedright\arraybackslash}p{0.16\linewidth}
  >{\raggedright\arraybackslash}X
  >{\raggedright\arraybackslash}X@{}}
\toprule
\thead{Split} & \thead{$n$ (lines; drugs; wells)} & \thead{model $\NIR$}
  & \thead{$\gap$, clustered on \emph{cell line}}
  & \thead{the same, clustered on \emph{drug}} \\
\midrule
\texttt{train} & $300$ ($39$; $76$; $136$) & $0.598$ $[0.549, 0.647]$
  & \ci{+0.0898}{+0.0382}{+0.1415} & \ci{+0.0898}{+0.0291}{+0.1506} \\
\addlinespace
\textbf{\texttt{unseen\_combo}} & \needsaudit{895} ($42$; $34$; $36$)
  & $0.533$ $[0.492, 0.575]$
  & \ci{+0.0332}{+0.0069}{+0.0594} & \ci{+0.0332}{-0.0031}{+0.0695} \\
\addlinespace
\texttt{unseen\_drug} & $199$ ($40$; $10$; $28$) & $0.459$ $[0.404, 0.514]$
  & \ci{-0.0315}{-0.0836}{+0.0206} & \ci{-0.0315}{-0.1049}{+0.0419} \\
\bottomrule
\end{tabularx}
\begin{tablenotes}[flushleft]\footnotesize
\item Every gap printed here is measured against the geometry-defined
\texttt{orth} partner, whose own score is empirically near chance on these
conditions; \texttt{opposite} figures never appear as a column. Model $\NIR$
intervals test against the chance value $0.5$ with no comparator. Intervals are
two-way cluster-robust over cell lines and wells, with a Student $t$ reference
on $\min(G_{\textrm{line}}, G_{\textrm{well}}) - 1$ degrees of freedom, which is
why the well counts are printed. The \texttt{unseen\_drug} arm spans $28$
wells, and $t_{27} = 2.05$ against $1.96$ widens its interval by five per cent
(\cref{sec:methods-data}).
\end{tablenotes}
\end{threeparttable}
\end{fullwidth}
\end{table}

\paragraph{The control behaves as the diagnosis predicts} The unseen-drug gap
against \texttt{opposite} is \ci{+0.0379}{-0.0083}{+0.0841}, positive where the
\texttt{orth}-partner estimate on the same conditions and predictions,
\ci{-0.0315}{-0.0836}{+0.0206}, is negative and covers zero as a control should.
That point estimate is the \texttt{unseen\_combo} gap with the sign reversed, so
the row is read as uninformative and not as a confirmed null. A cluster-robust
variance is asymptotic in the number of clusters, not conditions, and this arm
spans $40$ cell lines but only $28$ treatment wells, so its two-way variance is
carried on $27$ degrees of freedom. A cell-line-only cluster bootstrap puts the
\texttt{opposite} gap clear of zero at $[+0.0044, +0.0727]$; adding wells as the
second clustering axis returns it to zero, at the normal quantile as well as at
$t_{27} = 2.05$. What decides the row is the unit of clustering.

\paragraph{The held-out row settles prompt sensitivity only} On
\texttt{unseen\_combo} the gap clears zero clustered on cell line and well but
not clustered on drug and
well (\cref{fig:splits}a, right), so it survives only one of the two units the
claim needs. Drug is the axis
a claim about transferring a \emph{drug}-specific residual has to generalise
across, so
generalisation to an unseen treatment well of a seen drug is \notestablished.

\begin{figure}[htbp]
\centering
\includegraphics{fig-splits}
\caption[Chance-clearing on trained conditions only]{\textbf{The residual-frame
model clears chance on trained conditions only; the held-out-combination gap
excludes zero clustered on cell line but not clustered on drug, and the
unseen-drug gap flips sign.} Residual frame throughout, and $\NIR$ ranks each
prediction against the other drugs measured in the same cell line.
\textbf{(a)}~Left, model $\NIR$ against the chance value $0.5$ (solid rule), with
$95\%$ two-way cluster-robust intervals; $n = 300$, \needsaudit{895} and $199$ conditions on
\texttt{train}, \texttt{unseen\_combo} and \texttt{unseen\_drug}. The grey
dash on each row is that split's raw within-well split-half precision
reference, a measurement-precision figure computed from two halves of the same
well and not a replicate experiment; it carries no interval and is drawn as a
bare marker, and only \texttt{train} is filtered at the reliability line. Right,
on its own axis, the gap against the geometry-defined \texttt{orth} partner. Each
row has one point estimate, drawn twice with two intervals because the verdict
depends on which unit of generalisation is clustered: cell-line-clustered above,
drug-clustered below, each labelled with the Student~$t$ degrees of freedom it is
carried on. The grey band on the \texttt{unseen\_drug} row is that arm's
$80\%$-power detection limit around chance, $\pm 0.079$, the smallest departure
the arm could detect and not a confidence region. \textbf{(b)}~The per-condition
$\NIR$ distributions behind those three means, as reflected-boundary kernel
densities on a shared axis, each with its mean marked and the region at or below
chance shaded. The pooled $0.537$ is a mean over nearly flat distributions rather
than a location: even on \texttt{train}, $38.7\%$ of conditions sit at or below
chance, against $44.6\%$ and $55.8\%$ on the two held-out splits.}
\label{fig:splits}
\end{figure}

\begin{scopelimit}[carried into \cref{sec:res-lookup} and \cref{sec:limitations}]
The precision matters: \notestablished is not the same as \shownzero. The
drug-clustered interval reaches $+0.0695$, so an effect of up to roughly $+0.07$
remains compatible with these data at the
\needsaudit[arm size, two records disagree]{895} conditions this arm affords,
spread over $34$ drugs, a limit of power at the size the holdout gives,
reported as one. On \texttt{unseen\_drug}, $n = 199$ leaves a half-width of
$0.055$ around the model's own $\NIR$ and an $80\%$-power detection limit of
$0.079$ against chance, shaded on that row in \cref{fig:splits}a. Neither row
licenses a statement that the effect is absent.
\end{scopelimit}

% FLAGGED, NOT FIXED -- the two panel-B numbers below (train $0.955$ and the
% spread $0.006$) are inherited from a superseded reconstruction taken at a
% +0.109 reliability line. This document defines the line at +0.1086: see :214
% ("$+0.109$, resolving to $+0.1086$"), :222-223 ("none falls at or below the
% resolved $+0.1086$") and :986 ("499 conditions sit below the $+0.1086$ line").
% Recomputed from RESULTS_cluster/re_v3.json -> records[] at repro_cos > 0.1086:
%   train         kept 300  0.9555006 -> 0.956   (:222-223 requires all 300)
%   unseen_combo  kept 396  0.9605642 -> 0.961   (:986's 499-below requires 396)
%   unseen_drug   kept 111  0.9597734 -> 0.960
%   spread: raw 0.0050636 -> 0.005, and max(rounded) - min(rounded)
%           = 0.961 - 0.956 = 0.005. BOTH conventions give 0.005 at the resolved
%           line, so this is NOT the convention clash it was reported as. 0.006
%           is just the old line's rounded-endpoint spread, 0.961 - 0.955. At
%           0.109 the kept counts are 299 / 394 / 111, which contradict :222-223
%           and :986 directly.
% The $0.132$ in the same sentence is unaffected: panel A is read straight from
% CANONICAL_NUMBERS.json -> ceilings{}, and raw (0.13178) and rounded-endpoint
% (0.956 - 0.824) both give 0.132 there.
% NOT CHANGED HERE because the correction is a six-site edit and four sites are
% outside this file: 05-Limitations.tex:224 and :225, 03-Apparatus.tex:378, and
% figs/splitref.tex's spread node (:219) and caption (:244). That figure already
% DRAWS train as 0.956 at the resolved line while the prose here still prints
% 0.955, so the disagreement is live. Changing this site alone would replace a
% uniformly wrong number with an inconsistent document. Reported upward instead.
The three splits are not scored on comparable populations either, and the
difference is one of selection. Their raw within-well split-half precision
references, the grey dashes of \cref{fig:splits}a, are $0.956$, $0.824$ and
$0.836$, a spread of $0.132$; \texttt{train}
is filtered at a reliability line the two held-out splits are not, and imposing
that line on all three brings the references to $0.955$, $0.961$ and $0.960$, a
spread of $0.006$. The held-out shortfall therefore cannot be softened by appeal
to a lower noise floor.

% PLACEMENT ONLY. The other retraction in this section is guarded at the
% \needspace above; this one was not, and it broke 13 lines / 2 lines across the
% page. The two lines that landed overleaf are the ones that perform the
% withdrawal, and they arrived with no CORRECTION rule above them, so the reader
% met the retraction as unlabelled prose. \budgettick reserves 5\baselineskip,
% which is not enough: the box measures 247.1pt, 16.6 lines with its padding and
% frame, so 17 is the box entire and it either sets whole here or starts the
% next page whole.
% MEASURED COST, build of 2026-08-12: the guard carries ~2.6 lines of the
% chapter forward. That stops the closing claim of \cref{sec:res-lookup}
% breaking, and starts the coverage definition and the seen-drug-transfer claim
% breaking, all three in 04d-baselines.tex, each of which wants the same guard.
\needspace{17\baselineskip}%
\begin{withdrawn}{a training-exposure or memorisation advantage}
The difference between the trained and held-out estimates is
\ci{+0.0567}{-0.0014}{+0.1148} with cluster-robust $p = 0.0556$. An unclustered
permutation returns $p = 0.0142$, quoted only for contrast, because ignoring the
clustering this chapter mandates everywhere else flatters the result. It does
not clear the conventional threshold, and the two arms differ in three ways at
once in any case, in whether the condition was trained on, in which drugs they
contain, and in how their scored populations were selected. Restricting the
comparison to the $30$ drugs present in both arms reverses the sign, with the
trained arm sitting \emph{behind} the held-out one,
\ci{-0.0606}{-0.1506}{+0.0295}, $p = 0.1797$, on $79$ trained and $750$
held-out conditions, and the pooled difference is a composition effect: $83.8\%$
of held-out-combination conditions sit on those shared drugs against $26.3\%$
of trained ones. The claim of a training-exposure or memorisation advantage is
therefore withdrawn. None is established, the drug-matched estimate points the
other way, and the two arms are scored on populations selected by different
rules.
\end{withdrawn}

\paragraph{Two instruments no comparator defect can touch} First, the slope test.
Regressing each prediction's $\NIR$ on the cosine between the named drug's truth
and the target's truth gives \ci{+0.373}{+0.280}{+0.466} on \texttt{train},
\ci{+0.336}{+0.223}{+0.450} on \texttt{unseen\_combo} and
\ci{+0.342}{+0.215}{+0.469} on \texttt{unseen\_drug}, clustered on the drug. All
three exclude zero, and all three still exclude it clustered on the cell line
($[+0.272, +0.474]$, $[+0.241, +0.432]$ and $[+0.284, +0.400]$). Its points are
the three \emph{scrambled} partner prompts and nothing else, exactly three per
condition, so $900$, $\num{2685}$ and $597$ points on the three splits. On
\texttt{unseen\_drug} those partners are overwhelmingly drugs the model
\emph{was} trained on, so the slope shows that the model responds to a named
trained drug in a scoring context whose own drug was held out, and not that it
knows the held-out drug.

Second, the model's own score against chance, with no comparator at all. It
clears $0.5$ on \texttt{train}, by $+0.0981$
\nolinebreak[3]$[+0.0407, +0.1554]$ clustered on drug and well, and is not
distinguishable from chance on either held-out split
(\cref{tab:generalisation}).

\paragraph{A positive slope beside a chance-level score is not a contradiction}
The model has a coarse sense of what a named trained drug does, and not of what
it does in this well, and $\NIR$ asks the second question, ranking against the
$\approx 40$ other drugs in the same cell line. That is the drug main effect
without the part of the residual that is not shared across contexts, which
\cref{sec:res-transfer} measures at $44.7\%$ ($1 - \Tcoef$, with
$\Tcoef = 0.553$ \nolinebreak[3]$[0.514, 0.593]$), a figure that mixes cell line,
dose, well and estimator. It is also why the lookup wins in the next section:
\texttt{drug\_lookup} is that main effect, estimated straight from data far more
precisely than the model learned it. The condition-level picture agrees. The
cosine between prediction and own truth on trained conditions averages $+0.0389$
(sd $0.1119$, median $+0.0220$), spanning $[-0.3046, +0.3976]$ with $33.7\%$
clearing $|0.1|$.

\paragraph{One asymmetry attaches to every held-out number in this chapter} The
evaluation's truths are not refitted. The generic is reconstructed from the
build's own fit inventory and verified against the build's \texttt{fit\_digest},
and a mismatch aborts the run. The reliability line selects the \emph{training}
split, so \texttt{train} is by construction the reliability-surviving population,
and the two held-out splits are not filtered at all. In the unseen-combination
arm $\num{499}$ conditions sit below the $+0.1086$ line and $\num{74}$ have a
negative split-half cosine; in the unseen-drug arm, of $\num{199}$ conditions,
$\num{88}$ sit below the line and $\num{18}$ are negative. This is deliberate,
since filtering the evaluation set on its own outcome is the selection this
chapter removed, but it leaves the comparison conservative in the direction that
matters. The held-out arms are scored on their full, noisier populations while
the trained arm is not.

% PLACEMENT ONLY. This box is 23 lines and it was breaking after two of them,
% which is the one split `breakable' is not for: the accent rule, the words "The
% claim." and two lines at the foot of the page, and the whole verdict overleaf.
% It fits on one page with room to spare, so the reservation is the whole box.
\needspace{24\baselineskip}%
\begin{claim}
Re-encoding the target made the model respond to the combined drug-and-mechanism
prompt. That much is established and survives every stress: on trained
conditions it holds clustered on the cell line and on the drug, it holds against
chance with no comparator at all, and the comparator-free slope excludes zero on
all three splits under both clusterings (\ci{+0.373}{+0.280}{+0.466},
\ci{+0.336}{+0.223}{+0.450}, \ci{+0.342}{+0.215}{+0.469}).

Transfer to held-out combinations is \notestablished. The gap there is
$+0.0332$. Clustered on cell line and well it excludes zero,
\ci{+0.0332}{+0.0069}{+0.0594}; clustered on drug and well it does not,
\ci{+0.0332}{-0.0031}{+0.0695}. It is the drug axis that the claim has to
survive. The split is also weaker than its name: $30$ of its $34$ drugs and
$39$ of its $42$ cell lines appear in training, and $82\%$ of its conditions
recombine a seen drug with a seen cell line. The effect is a majority tilt
rather than a uniform one, positive for $20$ of $34$ drugs (sign test
$p = 0.39$).

What the slope on the \emph{unseen-drug} split does and does not add has to be
stated exactly. At \ci{+0.342}{+0.215}{+0.469} it shows that the prediction
still tracks the drug \emph{named in the prompt} when the condition's own drug was
held out. But the drugs named in that regression are the three scrambled
partners, which on this split are overwhelmingly training drugs, and the real
target-drug prompt contributes no point to it. The slope is evidence that the
model responds to named trained drugs in an unseen-drug scoring context, and not
evidence that it knows the held-out drugs.
\end{claim}
